\label{sec:compton-total-simplification}
\section{The total cross section}
\textbf{Roman~N.~Lee, Matthew~D.~Schwartz and Xiaoyuan Zhang have already obtained
an analytic expression for this total cross section~\cite{qed-processes-lsz}.
The expression developed here is an independent recalculation for
verification.}
The integrations of our virtual and real amplitudes are complete.
This chapter simplifies their analytic sum in a common energy-variable
basis. We begin with the sector assembly already derived, reduce its
terms to that basis, and collect their coefficients. No published
cross-section coefficient is used in this reduction.
We use
\begin{equation}
 X=\frac{s_{\mathrm{phys}}}{m^2}-1,\qquad
 q=\frac{\sqrt{X+4}-\sqrt X}{\sqrt{X+4}+\sqrt X},\qquad
 Y=\sqrt{\frac{X}{X+4}},\qquad X>0.
\end{equation}
The square roots are positive.  The Thomson coupling is denoted by
$\alpha_{\mathrm T}$, and $m$ is the electron mass.  Define
$s=1+X$ and $r_\omega=X/2$, so that the incident photon energy in the
electron rest frame is $m r_\omega$.  The complete result is
\begin{equation}
 \boxed{\sigma_{\mathrm{inc}}^{\mathrm{NLO}}(X)
 =\frac{\pi\alpha_{\mathrm T}^2}{m^2}B(X)
  +\frac{\alpha_{\mathrm T}^3}{m^2}C_1(X)
  +\frac1{m^2}O(\alpha_{\mathrm T}^4).}
 \label{eq:comp-simplified-total}
\end{equation}
Here the real contribution includes the complete two-photon final-state
phase space, with its identical-photon factor.  The remainder is at fixed
$X>0$.  The Born coefficient is
\begin{equation}
 \begin{aligned}
 B(X)=2\bigg[&\frac{1+r_\omega}{r_\omega^2}
 \left(\frac{2(1+r_\omega)}{1+2r_\omega}
       -\frac{\ln(1+2r_\omega)}{r_\omega}\right)\\
 &+\frac{\ln(1+2r_\omega)}{2r_\omega}
       -\frac{1+3r_\omega}{(1+2r_\omega)^2}\bigg].
 \end{aligned}
 \label{eq:comp-simplified-born}
\end{equation}
For the NLO coefficient, put
\begin{equation}
 \begin{aligned}
 \mathcal R_{\mathrm{rem}}(X)
 ={}&F_{\mathrm{rat}}^z+\Delta_{\mathrm{rat}}
       +F_{\log}^z+\Delta_{\log}+F_W^z\\
       &+F_D^z+\Delta_D+2F_{UW}^{\eta}.
 \end{aligned}
 \label{eq:comp-simplified-real-remainder}
\end{equation}
\begin{equation}
 \begin{aligned}
 C_1(X)={}&\widehat C_{V+\mathrm{UV}}(X)
 +\frac{\mathcal R_{\mathrm{rem}}(X)}{8X(1+X)}\\
 &-A(1+X)\ln\!\bigl(X\sqrt{1+X}\bigr)
 -\frac{\mathcal C_K(1+X)}{1+X}.
 \end{aligned}
 \label{eq:comp-simplified-master}
\end{equation}
All finite parts in $\mathcal R_{\mathrm{rem}}$ use the unnormalized
energy integrals of \eqref{eq:comp-assembly-real-sectors}.
Specifically, $F_j^z$ is the constant term of the endpoint expansion
with $z>\zeta$ after subtracting its power and logarithmic divergences.
The finite GPL endpoint formula is \eqref{eq:comp-real-gpl-endpoint-finite}.
The term $D$ restores the exact difference between $VW$ and the second
copy of $UW$, as defined in \eqref{eq:comp-real-bose-restoration}.
The shifts $\Delta_{\mathrm{rat}}$ and $\Delta_{\log}$ are given by
\eqref{eq:comp-real-simple-matching}.  The shift $\Delta_D$ includes the moving
boundary and double-pole terms in
\eqref{eq:compton-general-soft-matching}.  The $W$ sector has no such
shift.  The finite part $F_{UW}^{\eta}$ uses the analytic-regulator
definition and closed primitives of
\eqref{eq:comp-real-UW-eta-finite}--\eqref{eq:comp-real-UW-inner-closed}.
Thus $\mathcal R_{\mathrm{rem}}$ contains already integrated terms.  It
is not an omitted physical contribution or an unevaluated phase-space
integral.

The infrared coefficient $A(s)=\mathcal C_{\mathrm{IR}}(s)/s$ has
the integrated form \eqref{eq:comp-soft-integrated-ir}. The finite soft term
$\mathcal C_K(s)$ is given by \eqref{eq:comp-soft-K-angular-GPL}.
The three endpoint shifts are the finite regulator conversions
derived in \eqref{eq:compton-general-soft-matching}.
To define the first term in \eqref{eq:comp-simplified-master}, set
\begin{equation}
 C_{V+\mathrm{UV}}(X;\ell_\lambda)
 =C_V(1+X,\lambda_{\mathrm{phys}}/m)
 +\frac{F_{UV}^{\eta}(X)}{8X(1+X)}.
\end{equation}
Here $UV$ labels the real angular denominator pair, not an ultraviolet
counterterm. Separate the photon-mass logarithm as
\begin{equation}
 C_{V+\mathrm{UV}}(X;\ell_\lambda)
 =A(1+X)\ell_\lambda+\widehat C_{V+\mathrm{UV}}(X),
 \qquad \ell_\lambda=\ln(\lambda_{\mathrm{phys}}/m).
 \label{eq:comp-simplified-log-split}
\end{equation}
Its 23 coefficients are displayed below.  The hat denotes the constant
coefficient in this polynomial in $\ell_\lambda$, rather than a choice
of a large photon mass.  Equation~\eqref{eq:comp-simplified-master}
follows by canceling the photon-mass logarithm in
\eqref{eq:comp-assembly-common-regulator}.  It is independent of
$\lambda_{\mathrm{phys}}$.  The reductions below put every finite contribution into the same
alphabet before their coefficients are combined.

\section{The virtual contribution in a common function basis}
\label{sec:comp-virtual-energy-basis}
\subsection{A constant alphabet}
Use iterated integrals based at $X=0$, with the following kernels:
\begin{align}
 G(a,w;X)&=\int_0^X\omega_a(t)G(w;t),\qquad G(\varnothing;X)=1,\\
 \omega_a&=\frac{\dd X}{X-a}\quad(a=0,1,-1,-2,-4),\qquad
 \omega_r=\frac{\dd X}{\sqrt{X(X+4)}}.
\end{align}
The root is positive for $X>0$. Pure-zero words have the normalization
$G(0,\ldots,0;X)=\ln^nX/n!$. The letter $1$ occurs in combinations
whose apparent pole is canceled by an inner function vanishing at
$X=1$, as in $G(1,0;X)$ with kernel $\ln X/(X-1)$.

Define a new variable $q$, unrelated to electric charge, by
\begin{equation}
 q=\frac{\sqrt{X+4}-\sqrt X}{\sqrt{X+4}+\sqrt X},\qquad
 X=\frac{(1-q)^2}{q},\qquad 0<q<1.
\end{equation}
The physical square roots are positive.  Introduce
\begin{equation}
 \rho_\pm=\frac{1\pm\ii\sqrt3}{2},\qquad
 \gamma_\pm=\frac{3\pm\sqrt5}{2},\qquad
 \nu_a=\frac{\mathrm{d}q}{q-a}.
\end{equation}
Direct differentiation gives the complete pullback of the kernels:
\begin{align}
 \frac{\mathrm{d}X}{X}&=2\nu_1-\nu_0,\\
 \frac{\mathrm{d}X}{X+1}&=\nu_{\rho_+}+\nu_{\rho_-}-\nu_0,\\
 \frac{\mathrm{d}X}{X+2}&=\nu_{\ii}+\nu_{-\ii}-\nu_0,\\
 \frac{\mathrm{d}X}{X+4}&=2\nu_{-1}-\nu_0,\\
 \frac{\mathrm{d}X}{X-1}&=\nu_{\gamma_+}+\nu_{\gamma_-}-\nu_0,\\
 \frac{\mathrm{d}X}{\sqrt{X(X+4)}}&=-\nu_0.
 \label{eq:compton-equivalence-pullback}
\end{align}
Replacing each letter by its linear combination in these equations
converts every independently integrated word to constant-letter GPLs based at $q=1$.
The basepoint is approached from the left.  Its logarithmic convention
is fixed by the exact relation
\begin{equation}
 \ln X=2\ln(1-q)-\ln q.
\end{equation}
Writing $q=1-\delta$ gives $X=\delta^2/(1-\delta)$ and
$\ln X=2\ln\delta+O(\delta)$.  The leading coefficient of
$X$ in this relation is one, so it produces no finite logarithmic
constant.  This fixes words ending in the singular letter without an
undetermined additive constant.  All other boundary constants follow
by induction from the vanishing lower-weight words.  Complex letters
occur in conjugate pairs, with the same real path and conjugate boundary
values.  The resulting alphabet is
\begin{equation}
 \mathcal A=\{0,1,-1,\ii,-\ii,\rho_+,\rho_-,\gamma_+,\gamma_-\}.
\end{equation}
The transformed words have finite rational coefficient dictionaries.
All six kernel identities
in \eqref{eq:compton-equivalence-pullback} are checked symbolically.

\subsection{The vertex and self-energy contribution}
The entire independently integrated vertex and self-energy contribution
can already be reduced to constant letters.  In the direct channel, the
physical continuation gives
\begin{equation}
 \operatorname{Re}\operatorname{Li}_2(1+X+\ii0)
 =\frac{\pi^2}{6}-G(-1,0;X),\qquad
 \operatorname{Re}\ln(-X-\ii0)=\ln X.
\end{equation}
The imaginary part of the dilogarithm is fixed by its indicated boundary,
not by a fit to the final answer.  The interference coefficients in this
channel are real, so these real parts suffice for this contribution.
In the crossed channel the two angular endpoints are $y=X$ and
$y=X/(1+X)$.  At the latter endpoint, the substitution $y=t/(1+t)$ gives
\begin{equation}
 \frac{\mathrm{d}y}{y}=\frac{\mathrm{d}t}{t}
 -\frac{\mathrm{d}t}{t+1},\qquad
 \frac{\mathrm{d}y}{y-1}=-\frac{\mathrm{d}t}{t+1}.
\end{equation}
Its derivative at $t=0$ is one, so the logarithmic basepoint adds no
constant: $\ln y=\ln t-\ln(1+t)$.  Applying these replacements to the
independent angular primitive, including its rational coefficients,
gives exactly thirteen words:
\begin{equation}
 C_{V+\Sigma}(s,\lambda)
 =\sum_{w\in\mathcal W_V}v_w(X,\widehat c_Z)G(w;X),\qquad
 \widehat c_Z=-4-4\ln\lambda.
\end{equation}
Here $\lambda$ is measured in units of $m$. The symbol
$\widehat c_Z$ is the ultraviolet-finite field coefficient used in the
angular primitive, not the full counterterm $c_Z$ of Chapter 2.
With $G(\varnothing;X)=1$, the word set is
\begin{align}
 \mathcal W_V={}&\{\varnothing,(-1),(0),(-1,-1),(-1,0),(0,-1),(1,0),\\
 &(-1,-1,-1),(-1,-1,0),(-1,0,-1),(0,-1,-1),(0,-1,0),(0,1,0)\}.
\end{align}
Every $v_w$ is rational in $X$ and $\widehat c_Z$.  All explicit $\pi^2$ constants
cancel in this representation.  For example,
\begin{align}
 v_{(0,1,0)}&=v_{(0,-1,0)}=\frac{4(X+1)}{X^3},\\
 v_{(-1,-1,0)}&=-\frac{4(X-1)}{X^3},\qquad
 v_{(1,0)}=\frac{X^2-9X+16}{X^3}.
\end{align}
The full coefficient and derivative dictionaries are obtained from the
independent trace and angular-primitive files, without using the
published total.  This reduces the complete vertex and self-energy subset. It is not
separately gauge invariant and must be combined with the spanning boxes.

\subsection{The direct spanning box}
The direct spanning-box channel can also be reduced completely.  In its
angular integral let
\begin{equation}
 b=\sqrt{\frac{X}{X+4}},\qquad
 y=\frac{X-(X+4)v^2}{1-v^2},\qquad
 0<v<V=\frac{X}{X+2}.
\end{equation}
Its input GPL words use only $0,\pm1$.  Exact rational integration in
$v$ adds first letters $\pm b$, with no other roots.  Setting
$X=4b^2/(1-b^2)$ makes every rational coefficient and partial-fraction
pole rational in $b$.  The resulting angular primitive has been
checked by differentiation, with every GPL coefficient canceling
identically for arbitrary $b$.

At the physical endpoint the variable $q$ gives
\begin{equation}
 b=\frac{1-q}{1+q},\qquad V=\frac{(1-q)^2}{1+q^2}.
\end{equation}
Both quantities are rational in $q$.  Let $\mathcal G$ denote GPLs
based at $q=1$, with the same real left tangent used above.  The
weight-one transformations are
\begin{align}
 G(0;V)&=2\mathcal G(1;q)-\mathcal G(\ii;q)
                   -\mathcal G(-\ii;q)-\ln2,\\
 G(1;V)&=\mathcal G(0;q)-\mathcal G(\ii;q)-\mathcal G(-\ii;q),\\
 G(-1;V)&=\mathcal G(\rho_+;q)+\mathcal G(\rho_-;q)
                  -\mathcal G(\ii;q)-\mathcal G(-\ii;q),\\
 G(b;V)&=2\mathcal G(0;q)-\mathcal G(\ii;q)-\mathcal G(-\ii;q),\\
 G(-b;V)&=-\mathcal G(\ii;q)-\mathcal G(-\ii;q).
\end{align}
These follow by taking logarithms of the corresponding rational ratios,
with their constants fixed at $q=1$.  In particular, for
$q=1-\delta$ one has $V=\delta^2/2+O(\delta^3)$, so the $-\ln2$
in the first line is necessary.

Higher weights follow recursively from the total differential of a
moving-letter GPL.  All differences between $0,\pm1,\pm b$ and $V$
are rational functions of $q$.  Their logarithmic derivatives have
simple poles only in
\begin{equation}
 \mathcal A_s=\{0,1,-1,\ii,-\ii,\rho_+,\rho_-\}.
\end{equation}
After a lower-weight word has been converted, integration of each term
$\nu_a\mathcal G(w;q)$ prepends the letter $a$.  Every original word
containing a nonzero letter vanishes as $q\to1$, since $V=O(\delta^2)$
while $b=O(\delta)$.  Thus its integration constant is zero.  The
all-zero words are fixed by powers of the first weight-one identity.
Coincident adjacent letters cancel in the total differential before
individual logarithmic derivatives are taken.  The transformation
checks these cancellations exactly, including the cancellation of
any spurious pure powers of $\ln(1-q)$ at the boundary.

Finally, the scalar factors multiplying the angular words satisfy
\begin{equation}
 d_s=\frac{G(-1,0;X)}{X},\qquad \ell_X=G(0;X).
\end{equation}
Transforming these factors with the kernel map above and applying
shuffle identities completes the direct-box result as 101 constant-letter
words of weight at most three.  Its rational coefficients depend only
on $q$ and $\ln\lambda$, with possible constant $\ln2$ factors.  The
normalization is $C_{\mathrm{box},s}=J_s/(2X^2)$.
This reduction uses only the independent box coefficients.

\subsection{The crossed spanning box in the energy basis}
Section~\ref{sec:comp-crossed-first} has already integrated the rational
crossed subgroup in $y$, before introducing the angular square root.
Its endpoint maps contain only the letters $-1,0,1$. Applying the
kernel pullbacks above converts those words to the same $q$ basis.
The remaining angular integral has no auxiliary $\rho_u$ letter.

\subsection{The complete virtual contribution}
After the crossed subgroup is integrated as just described, its remaining
seventeen angular words involve only $0,\pm1,\pm b$.  The same total
differential construction used for the direct box therefore applies.
Combining it with the crossed subgroup gives 110 constant-letter words
for $J_u$.  The two independent constructions of the subgroup endpoint
coefficients agree as exact rational functions.  The remaining angular
primitive has also been checked by an exact derivative, coefficient by
coefficient, for arbitrary kinematics.

The entire independently computed virtual contribution is consequently
in the comparison basis:
\begin{equation}
 C_V(q,\lambda)=\sum_{w\in\mathcal W}V_w(q,\ln\lambda)\mathcal G(w;q).
\end{equation}
There are 116 nonzero words.  All explicit $\ln2$ coefficients introduced
by intermediate changes of endpoint cancel in this complete virtual
sum.  This is an algebraic cancellation of coefficients, not a numerical
inference.  The coefficients still depend on $\ln\lambda$, as expected
for the virtual contribution.

\subsection{Exact inversion of the kernel transformation}
The common basis can be reduced without choosing functions from a
published answer. Let $M$ be the $9\times6$ matrix of coefficients
in the six kernel pullbacks above. A row labels a letter in
$\mathcal A$, and a column labels one of the six $X$ kernels. Its
six columns are linearly independent. Select six independent rows
with a projection matrix $P$ and define
\begin{equation}
 L=(PM)^{-1}P,\qquad LM=I_6.
\end{equation}
For weight $n$, let $v_n$ denote the entire coefficient vector of
the independently integrated $q$ words, with zero entries for absent
words. Applying the tensor product of $n$ copies of $L$ gives the
candidate $X$ coefficients $c_n$. The necessary and sufficient
reconstruction test is
\begin{equation}
 c_n=L^{\otimes n}v_n,\qquad
 M^{\otimes n}c_n-v_n=0,\qquad 0\leq n\leq3.
\end{equation}
Every component of the second equation is checked as an exact
rational identity, including the rows not selected by $P$. Thus the
choice of independent rows cannot hide an omitted word. The
weight-zero equation includes constants. The common basepoint and
logarithmic normalization established above make this an equality
of integrals, including their boundary constants.

The remaining coefficient functions are rational in $q$. Polynomial
reduction using $q^2-(X+2)q+1=0$ reduces each one to $A(X)+qB(X)$.
The physical root is fixed by $0<q<1$:
\begin{equation}
 q=\frac{X+2-\sqrt{X(X+4)}}2.
\end{equation}
The other root exceeds one for $X>0$. Substitution therefore gives
coefficients rational in $X$ and $Y=\sqrt{X/(X+4)}$, with no branch
choice left to numerical evaluation.


\section{The virtual and UV contribution}
\paragraph{Common functions and exact alphabet reduction.}
The physical kernels required for the UV term are
\begin{equation}
 \omega_0=\frac{\dd X}{X},\qquad
 \omega_{-1}=\frac{\dd X}{X+1},\qquad
 \omega_{-4}=\frac{\dd X}{X+4},\qquad
 \omega_r=\frac{\dd X}{\sqrt{X(X+4)}}.
\end{equation}
Their exact pullbacks provide a finite algebraic test of whether the
apparently larger endpoint alphabet is needed:
\begin{align}
 \omega_0&=2\frac{\dd q}{q-1}-\frac{\dd q}{q},&
 \omega_{-4}&=2\frac{\dd q}{q+1}-\frac{\dd q}{q},\\
 \omega_{-1}&=\frac{\dd q}{q-\rho_+}+\frac{\dd q}{q-\rho_-}
                    -\frac{\dd q}{q},&
 \omega_r&=-\frac{\dd q}{q},\qquad
 \rho_\pm=\frac{1\pm \ii\sqrt3}{2}.
\end{align}
The endpoint $X=0$ corresponds to the same left tangential base point
$q=1$ used in the preceding calculation.  In particular,
$\ln X=2\ln(1-q)-\ln q$ has no additional constant there.

For each weight through three, form every tensor word in these four
kernels and expand it with the displayed pullbacks.  This gives a matrix
with integer entries.  Solving its coefficient equations against the
independently integrated UV dictionary, and substituting the solution
back into every equation, proves an identity of iterated integrals with
the same base point.  No comparison cross section enters this operation.
The rational prefactors are simplified using
\begin{equation}
 q^2-(X+2)q+1=0,\qquad
 q=\frac{X+2-(X+4)Y}{2}.
\end{equation}
Thus they belong to $\mathbb Q(X)+Y\mathbb Q(X)$, without auxiliary roots.
The two UV dictionaries each satisfy all 156 coefficient equations of
this transformation.  Their normalization remains that of the original
real integral: the contribution to the coefficient of $\alpha_{\mathrm T}^3/m^2$
is obtained by multiplying the UV finite part by $1/[8X(1+X)]$.

Adding this transformed finite part to the full virtual contribution
requires only the same 23 $X$ words already obtained for the virtual
amplitude.  The files accompanying this calculation give the separate
UV coefficients and the summed coefficients, so the normalization and
every cancellation can be checked before other contributions are added.
The 23-word sum by itself is not infrared finite.  The other terms appear with their complete endpoint definitions in
\eqref{eq:comp-simplified-master}.  The following sections reduce them to the same small alphabet,
so that their coefficients can be added algebraically.

After exact removal of vanishing coefficients, the UV finite dictionary
contains twelve nonzero $X$ words (rather than its original 41 $q$ words),
and its pole dictionary contains three.  All twelve finite words already
occur in the full virtual dictionary.  Thus their sum still contains
23 words.  Its coefficients are displayed explicitly below.
\paragraph{Explicit partial sum.}
Write $G_{a_1,\ldots,a_n}=G(a_1,\ldots,a_n;X)$ and $G_{\varnothing}=1$,
with the kernels above, $\omega_{-2}=\dd X/(X+2)$ and $\omega_1=\dd X/(X-1)$.
These iterated integrals have base point zero:
\begin{equation}
 G_{a,w}(X)=\int_0^X\omega_a(t)G_w(t),\qquad G_{\underbrace{\scriptstyle 0,\ldots,0}_{n}}(X)=\frac{\ln^nX}{n!}.
\end{equation}
The words ending in zero that contain the letter $1$ are understood by their continuous real extension at $X=1$.
Let $\ell_\lambda=\ln(\lambda_{\mathrm{phys}}/m)$.  The following expression includes the complete virtual contribution and the finite UV real contribution:
\begin{equation}
 C_{V+\mathrm{UV}}(X;\ell_\lambda)=\sum_{w\in\mathcal W}c_w(X;\ell_\lambda)G_w(X).
\end{equation}
The set $\mathcal W$ and every coefficient are specified by the following equations.
\begin{equation}
\begin{aligned}
144 X^{4} \left(X + 1\right)^{2}\,c_{\varnothing}&={}\bigl(- 5 X^{9} - 2 X^{8} + 141 X^{7}\bigr)\\
&\quad{}+\bigl(486 X^{6} - 720 X^{5} \ell_\lambda + 1380 X^{5}\bigr)\\
&\quad{}+\bigl(- 8640 X^{4} \ell_\lambda + 2128 X^{4} - 14400 X^{3} \ell_\lambda\bigr)\\
&\quad{}+\bigl(9392 X^{3} - 6912 X^{2} \ell_\lambda + 15616 X^{2}\bigr)\\
&\quad{}+\bigl(10752 X + 3072\bigr).
\end{aligned}
\end{equation}
\begin{equation}
\begin{aligned}
24 X^{5} \left(X + 1\right)^{2} \left(X + 4\right)\,c_{-1}&={}\bigl(- 5 X^{9} - 32 X^{8} - 48 X^{7} \ell_\lambda\bigr)\\
&\quad{}+\bigl(29 X^{7} + 1056 X^{6} \ell_\lambda + 80 X^{6}\bigr)\\
&\quad{}+\bigl(8832 X^{5} \ell_\lambda - 1656 X^{5} - 7504 X^{4}\bigr)\\
&\quad{}+\bigl(21312 X^{4} \ell_\lambda + 20352 X^{3} \ell_\lambda - 17088 X^{3}\bigr)\\
&\quad{}+\bigl(6912 X^{2} \ell_\lambda - 18944 X^{2} - 9728 X\bigr)\\
&\quad{}+\bigl(-2048\bigr).
\end{aligned}
\end{equation}
\begin{equation}
\begin{aligned}
36 X^{4} \left(X - 1\right) \left(X + 1\right)^{3}\,c_{0}&={}\bigl(3 X^{11} + 18 X^{10} - 42 X^{9}\bigr)\\
&\quad{}+\bigl(- 12 X^{8} + 699 X^{7} + 1336 X^{6}\bigr)\\
&\quad{}+\bigl(- 3694 X^{5} - 10404 X^{4} - 3680 X^{3}\bigr)\\
&\quad{}+\bigl(7616 X^{2} + 6912 X + 1536\bigr).
\end{aligned}
\end{equation}
\begin{equation}
\begin{aligned}
12 X^{5} \left(X + 1\right) \left(X + 4\right)\,c_{-1,-1}&={}\bigl(X^{8} - 51 X^{7} - 48 X^{6} \ell_\lambda\bigr)\\
&\quad{}+\bigl(- 596 X^{6} - 48 X^{5} \ell_\lambda - 900 X^{5}\bigr)\\
&\quad{}+\bigl(1056 X^{4} \ell_\lambda + 3772 X^{4} + 2208 X^{3} \ell_\lambda\bigr)\\
&\quad{}+\bigl(8448 X^{3} + 1152 X^{2} \ell_\lambda + 5824 X^{2}\bigr)\\
&\quad{}+\bigl(2816 X + 1024\bigr).
\end{aligned}
\end{equation}
\begin{equation}
\begin{aligned}
6 X^{5} \left(X + 1\right)^{2} \left(X + 4\right)\,c_{-1,0}&={}\bigl(X^{9} - 2 X^{8} - 77 X^{7}\bigr)\\
&\quad{}+\bigl(- 410 X^{6} - 650 X^{5} + 1558 X^{4}\bigr)\\
&\quad{}+\bigl(7372 X^{3} + 9744 X^{2} + 5376 X\bigr)\\
&\quad{}+\bigl(1024\bigr).
\end{aligned}
\end{equation}
\begin{equation}
\begin{aligned}
X^{3} \left(X + 1\right)^{2} \left(X + 4\right)\,c_{-2,-1}&={}\bigl(- 2 X^{5} - 52 X^{4} - 292 X^{3}\bigr)\\
&\quad{}+\bigl(- 632 X^{2} - 576 X - 192\bigr).
\end{aligned}
\end{equation}
\begin{equation}
\begin{aligned}
6 X^{5} \left(X + 1\right)^{2} \left(X + 4\right)\,c_{0,-1}&={}\bigl(X^{9} + 10 X^{8} + 7 X^{7}\bigr)\\
&\quad{}+\bigl(4 X^{6} + 1162 X^{5} + 5536 X^{4}\bigr)\\
&\quad{}+\bigl(11008 X^{3} + 10944 X^{2} + 5376 X\bigr)\\
&\quad{}+\bigl(1024\bigr).
\end{aligned}
\end{equation}
\begin{equation}
\begin{aligned}
X^{3}\,c_{1,0}&={}\bigl(- X^{2} + X - 2\bigr).
\end{aligned}
\end{equation}
\begin{equation}
\begin{aligned}
X^{4} \left(X + 4\right)\,c_{r,-1}&={}\bigl(3 X^{5} Y + 12 X^{4} Y \ell_\lambda + 36 X^{4} Y\bigr)\\
&\quad{}+\bigl(24 X^{3} Y \ell_\lambda + 104 X^{3} Y - 120 X^{2} Y\bigr)\\
&\quad{}+\bigl(- 288 X^{2} Y \ell_\lambda - 960 X Y \ell_\lambda - 768 X Y\bigr)\\
&\quad{}+\bigl(- 576 Y \ell_\lambda - 256 Y\bigr).
\end{aligned}
\end{equation}
\begin{equation}
\begin{aligned}
X^{3}\,c_{-1,-1,-1}&={}\bigl(2 X^{2} + 16 X + 22\bigr).
\end{aligned}
\end{equation}
\begin{equation}
\begin{aligned}
X^{3}\,c_{-1,-1,0}&={}\bigl(12 X^{2} - 20 X - 48\bigr).
\end{aligned}
\end{equation}
\begin{equation}
\begin{aligned}
X^{3}\,c_{-1,-2,-1}&={}\bigl(4 X^{2} - 16 X - 24\bigr).
\end{aligned}
\end{equation}
\begin{equation}
\begin{aligned}
X^{3}\,c_{-1,0,-1}&={}\bigl(4 X^{2} - 4 X - 12\bigr).
\end{aligned}
\end{equation}
\begin{equation}
\begin{aligned}
X^{4} \left(X + 4\right)\,c_{-4,r,-1}&={}\bigl(- 18 X^{4} Y - 48 X^{3} Y + 360 X^{2} Y\bigr)\\
&\quad{}+\bigl(1296 X Y + 672 Y\bigr).
\end{aligned}
\end{equation}
\begin{equation}
\begin{aligned}
X^{3}\,c_{0,-1,-1}&={}\bigl(- X^{2} - 8 X + 8\bigr).
\end{aligned}
\end{equation}
\begin{equation}
\begin{aligned}
X^{3}\,c_{0,-1,0}&={}\bigl(2 X^{2} + 4 X - 12\bigr).
\end{aligned}
\end{equation}
\begin{equation}
\begin{aligned}
X^{3}\,c_{0,1,0}&={}\bigl(2 X^{2} + 4 X - 12\bigr).
\end{aligned}
\end{equation}
\begin{equation}
\begin{aligned}
X^{4} \left(X + 4\right)\,c_{0,r,-1}&={}\bigl(3 X^{4} Y - 6 X^{3} Y - 144 X^{2} Y\bigr)\\
&\quad{}+\bigl(- 384 X Y - 480 Y\bigr).
\end{aligned}
\end{equation}
\begin{equation}
\begin{aligned}
X^{4} \left(X + 4\right)\,c_{r,-1,-1}&={}\bigl(21 X^{4} Y + 54 X^{3} Y - 432 X^{2} Y\bigr)\\
&\quad{}+\bigl(- 1536 X Y - 768 Y\bigr).
\end{aligned}
\end{equation}
\begin{equation}
\begin{aligned}
X^{4} \left(X + 4\right)\,c_{r,-1,0}&={}\bigl(- 6 X^{4} Y - 36 X^{3} Y + 192 X Y\bigr)\\
&\quad{}+\bigl(- 192 Y\bigr).
\end{aligned}
\end{equation}
\begin{equation}
\begin{aligned}
X^{4} \left(X + 4\right)\,c_{r,-2,-1}&={}\bigl(- 8 X^{4} Y - 16 X^{3} Y + 192 X^{2} Y\bigr)\\
&\quad{}+\bigl(640 X Y + 384 Y\bigr).
\end{aligned}
\end{equation}
\begin{equation}
\begin{aligned}
X^{4} \left(X + 4\right)\,c_{r,0,-1}&={}\bigl(6 X^{4} Y + 4 X^{3} Y - 192 X^{2} Y\bigr)\\
&\quad{}+\bigl(- 576 X Y - 576 Y\bigr).
\end{aligned}
\end{equation}
\begin{equation}
\begin{aligned}
X^{3}\,c_{r,r,-1}&={}\bigl(15 X^{2} - 12 X - 96\bigr).
\end{aligned}
\end{equation}

\section{A shorter infrared coefficient}
\label{sec:comp-ir-simplification}
The two roots in \eqref{eq:comp-soft-integrated-ir} are reciprocal.
This relation reduces its endpoint sum to one dilogarithmic function.
Set
\begin{equation}
 X=s-1>0,\qquad Y=\sqrt{\frac{X}{X+4}},\qquad
 q=\frac{X+2-\sqrt{X(X+4)}}2,\qquad L=\ln(1+X).
\end{equation}
Then $0<q<1$, $X=(1-q)^2/q$, and $z_-=q$, $z_+=q^{-1}$.
Define
\begin{equation}
 \mathcal G(X)=\frac12\ln^2q
 +\frac13\operatorname{Li}_2(-q^3)
 -\operatorname{Li}_2(-q)-\frac{\pi^2}{18}.
 \label{eq:comp-ir-short-function}
\end{equation}
Both dilogarithm arguments lie between $-1$ and zero.  All functions
in this expression are consequently real on the physical interval,
with their branches fixed without a continuation across a cut.
The integrated infrared coefficient is
\begin{equation}
 \begin{aligned}
 A(1+X)={}&-\frac{2(X^2-4X-6)}{X^3}L^2
 -\frac{2P(X)}{X^3(1+X)^2(X+4)}L\\
 &+\frac{12YQ(X)}{X^4(X+4)}\mathcal G(X)
 -\frac{R(X)}{X^2(1+X)^2},
 \end{aligned}
 \label{eq:comp-ir-short}
\end{equation}
where
\begin{align}
 P(X)&=X^5-22X^4-184X^3-444X^2-424X-144,\\
 Q(X)&=X^4+2X^3-24X^2-80X-48,\\
 R(X)&=5X^3+60X^2+100X+48.
\end{align}
The normalization is $A(s)=\mathcal C_{\mathrm{IR}}(s)/s$, so
$A(s)\ln(\lambda_{\mathrm{phys}}/m)$ is the logarithmic term of the
complete virtual coefficient in units of $\alpha_{\mathrm T}^3/m^2$.
Equation~\eqref{eq:comp-ir-short} is exact for every $X>0$.
Its apparent threshold poles cancel in the combined expression.

To derive it directly from the angular integral, denote the endpoint
difference of a primitive by $\Delta J=J(s)-J(1)$.  The two logarithmic
pole primitives in \eqref{eq:comp-soft-integrated-ir} satisfy
\begin{align}
 \Delta H_1(q)&=-\ln q, & \Delta H_1(q^{-1})&=\ln q,\\
 \Delta H_2(q)&=1, & \Delta H_2(q^{-1})&=1.
\end{align}
Their residues obey $c_{-,q}=c_{-,q^{-1}}$ and
$d_{-,q}+d_{-,q^{-1}}=8$.  Thus the logarithms of the negative-sign
sector cancel, while its double poles give the constant eight.
For the positive-sign sector, $c_{+,q}=-c_{+,q^{-1}}$.
Use the real dilogarithm inversion identity
\begin{equation}
 \operatorname{Re}\operatorname{Li}_2(u)
 +\operatorname{Re}\operatorname{Li}_2(u^{-1})
 =\frac{\pi^2}{3}-\frac12\ln^2u,\qquad u>0,
\end{equation}
to combine its simple-pole endpoint functions into
\begin{equation}
 \begin{aligned}
 \Delta J_1(q)-\Delta J_1(q^{-1})
 ={}&2\operatorname{Li}_2(q)-\operatorname{Li}_2(q/s)
  -\operatorname{Li}_2(qs)\\
 &-2L\ln q-\frac12L^2\equiv\mathcal F(q).
 \end{aligned}
\end{equation}
Here $qs=1-q+q^2<1$ and $q/s=q^2/(1-q+q^2)<1$.
Differentiation and the identity
$(1+q)(1-q+q^2)=1+q^3$ give
\begin{equation}
 \frac{\dd\mathcal F}{\dd q}=-\frac{3L}{q},\qquad
 \frac{\dd\mathcal G}{\dd q}=-\frac{L}{q},\qquad
 \mathcal F(1)=\mathcal G(0)=0.
\end{equation}
In the last boundary value the argument of $\mathcal G$ is $X$,
which tends to zero as $q\to1$.  Therefore
$\mathcal F(q)=3\mathcal G(X(q))$, with no undetermined constant.
The double-pole primitives supply only elementary terms,
\begin{align}
 \Delta J_2(q)&=-\frac{L}{s-q}-\frac{L+\ln q}{q},\\
 \Delta J_2(q^{-1})&=-\frac{L}{s-q^{-1}}+q(\ln q-L).
\end{align}
Substituting these differences and the elementary power primitives
into \eqref{eq:comp-soft-integrated-ir} yields
\eqref{eq:comp-ir-short}.  This derivation uses only the original
Born-weighted soft angular integral.

An independent internal check differentiates the complete virtual
23-word result with respect to $\ln(\lambda_{\mathrm{phys}}/m)$.
Its four nonzero coefficients agree algebraically with
\eqref{eq:comp-ir-short}, using
$G(-1;X)=L$, $G(-1,-1;X)=L^2/2$ and
$G(r,-1;X)=\mathcal G(X)$.  This check does not set any integration
constant in the soft calculation.

\section{Reduction of the virtual root functions}
\subsection{Collecting the coefficients}
A compact expression need not eliminate the $G$ functions. It is useful
first to combine their coefficients, before expanding any function into
classical polylogarithms. For the six root words in the preceding table,
the shuffle product gives
\begin{equation}
 G_{0,r,-1}=\ln X\,G_{r,-1}-G_{r,0,-1}-G_{r,-1,0}.
 \label{eq:comp-root-compact-shuffle}
\end{equation}
Thus $G_{0,r,-1}$ is redundant. After this substitution, the coefficients
of $G_{r,0,-1}$ and $G_{r,-1,0}$ are proportional with ratio $1:-3$.
Define
\begin{align}
 p_F(X)&=7X^4+18X^3-144X^2-512X-256,\\
 p_K(X)&=3X^4+10X^3-48X^2-192X-96,\\
 p_W(X)&=3X^4+8X^3-60X^2-216X-112,\\
 p_Z(X)&=X^4+2X^3-24X^2-80X-48,\\
 p_0(X)&=X^4-2X^3-48X^2-128X-160.
 \label{eq:comp-root-compact-polynomials}
\end{align}
The contribution $C_{\mathrm{six}}$ of these six original words to
$\widehat C_{V+\mathrm{UV}}$ becomes
\begin{equation}
 \begin{aligned}
 C_{\mathrm{six}}(X)=\frac{Y}{X^4(X+4)}\bigl\{&
 3p_FG_{r,-1,-1}+p_K(G_{r,0,-1}-3G_{r,-1,0})\\
 &-6p_WG_{-4,r,-1}-8p_ZG_{r,-2,-1}
       +3p_0\ln X\,G_{r,-1}\bigr\}.
 \end{aligned}
 \label{eq:comp-root-compact-G}
\end{equation}
All coefficients are obtained by exact rational collection of the original
table. This removes one weight-three function and gives the other five a
common denominator. The $G$ notation keeps this relation shorter than its
fully expanded classical form. The identities below provide alternative
representations for evaluation and further algebraic reductions.

\subsection{Auxiliary classical identities}
The rational-letter terms of the preceding coefficient list do not
require separate GPL evaluations. Put $L=\ln(1+X)$, $M=\ln X$ and
$v=1/(1+X)$, with $X>0$. The elementary reductions are
\begin{align}
 G_{-1}&=L,\qquad G_0=M,\qquad
 G_{-1,-1}=\tfrac12L^2,\qquad G_{-1,-1,-1}=\tfrac16L^3,\\
 G_{0,-1}&=-\operatorname{Li}_2(-X),\qquad
 G_{-1,0}=LM+\operatorname{Li}_2(-X),\\
 G_{1,0}&=\frac{\pi^2}{6}-\operatorname{Li}_2(1-X),\\
 G_{0,-1,0}&=-M\operatorname{Li}_2(-X)+2\operatorname{Li}_3(-X),\\
 G_{0,1,0}&=-M\mathcal L_2(X)+2\mathcal L_3(X).
\end{align}
Here $\mathcal L_n(X)=\operatorname{Re}\operatorname{Li}_n(X+\ii0)$.
The real boundary value is fixed by integrating
$\dd\mathcal L_2/\dd X=-\ln|1-X|/X$ and
$\dd\mathcal L_3/\dd X=\mathcal L_2/X$ from the interval $0<X<1$.
It is not selected by fitting numerical values. The combination
for $G_{1,0}$ is regular at $X=1$.

One additional classical combination suffices for the three words
with two letters $-1$ and one zero:
\begin{align}
 T(X)&=\tfrac12L^2M-\tfrac13L^3
       -L\operatorname{Li}_2(v)-\operatorname{Li}_3(v)+\zeta(3),\\
 G_{0,-1,-1}&=T(X),\\
 G_{-1,0,-1}&=-L\operatorname{Li}_2(-X)-2T(X),\\
 G_{-1,-1,0}&=\tfrac12L^2M+L\operatorname{Li}_2(-X)+T(X).
\end{align}
For example, substituting $v=1/(1+X)$ into
$\tfrac12\int_0^X\ln^2(1+t)\,\dd t/t$ gives the displayed $T$.
Its derivative is $L^2/(2X)$ and its lower limit is zero.
The last two lines then follow from the shuffle product.
All dilogarithms and trilogarithms with argument $v$ are on
$0<v<1$. These substitutions reduce the number of nonclassical
words in the partial sum without changing any rational coefficient
or the photon-mass convention.

The repeated square-root kernel can also be integrated classically.
Use the same $q$ as in \eqref{eq:comp-ir-short-function}. Since
$\omega_r=-\dd q/q$, integrating that independently derived expression
once more gives
\begin{equation}
 G_{r,r,-1}=\operatorname{Li}_3(-q)
 -\frac19\operatorname{Li}_3(-q^3)-\frac16\ln^3q
 +\frac{\pi^2}{18}\ln q+\frac23\zeta(3).
\end{equation}
Differentiating with respect to $q$ gives
$-\mathcal G(X(q))/q$. At $q=1$ the expression vanishes, using
$\operatorname{Li}_3(-1)=-3\zeta(3)/4$, which fixes the constant.
Its two trilogarithms have negative real arguments.
The remaining rational-letter word also has a short classical form.
With $u=1+X$, its inner weight-two function is
$\operatorname{Li}_2(-u)+\ln u\ln(1+u)+\pi^2/12$.
Integrating it with $\dd u/u$ gives
\begin{equation}
 G_{-1,-2,-1}=2\operatorname{Li}_3(-1-X)
 -L\operatorname{Li}_2(-1-X)+\frac{\pi^2}{12}L
 +\frac32\zeta(3).
\end{equation}
The last constant follows from the zero value at $X=0$.
Together with \eqref{eq:comp-ir-short-function}, these reductions
leave only six nonclassical words in this partial sum:
\begin{equation}
 \left\{G_{-4,r,-1},G_{0,r,-1},G_{r,-1,-1},
 G_{r,-1,0},G_{r,-2,-1},G_{r,0,-1}\right\}.
\end{equation}

\paragraph{A classical primitive for the squared logarithm.}
The word $G_{r,-1,-1}$ can be removed from this list without
introducing another iterated integral. Define
\begin{align}
 a&=e^{\ii\pi/3},\quad b=\overline a,\quad X=(1-q)^2/q,\quad 0<q<1,\\
 A&=1-aq,\quad B=1-bq,\quad R=A/B.
\end{align}
The quadratic factorization gives $AB=1-q+q^2$ and hence
$\ln(1+X)=\log A+\log B-\ln q$.
Here and below $\log$ is the principal complex logarithm.
Introduce the classical function
\begin{equation}
 \mathcal J(z,\rho)=\log^2z\log(1-z/\rho)
 +2\log z\operatorname{Li}_2(z/\rho)
 -2\operatorname{Li}_3(z/\rho).
\end{equation}
Differentiation of the polylogarithms cancels the terms proportional
to $\operatorname{Li}_2$ and $\log(1-z/\rho)$, leaving
\begin{equation}
 \frac{\partial\mathcal J(z,\rho)}{\partial z}
 =\frac{\log^2 z}{z-\rho}.
\end{equation}
For the ratio $R$, the inverse transformation is
$q=(1-R)/(a-bR)$. Consequently
\begin{equation}
 \frac{\dd q}{q}=\frac{\dd R}{R-1}
                   -\frac{\dd R}{R-a/b}.
\end{equation}
This second pole is essential. Applying this identity to
$2\log A\log B=\log^2A+\log^2B-\log^2R$ reduces the
product of logarithms to the same primitive. The remaining mixed
term follows from
\begin{equation}
 \frac{\dd}{\dd q}\left[-\ln q\operatorname{Li}_2(aq)
                  +\operatorname{Li}_3(aq)\right]
 =\frac{\ln q\log(1-aq)}q.
\end{equation}
Combining these operations gives
\begin{align}
 \mathcal F(q)&=2\mathcal J(A,1)+2\mathcal J(B,1)
                  -\mathcal J(R,1)+\mathcal J(R,a/b)
                  +\mathcal D(q),\\
 \mathcal D(q)&=2\ln q\left[\operatorname{Li}_2(aq)
                         +\operatorname{Li}_2(bq)\right]
 -2\left[\operatorname{Li}_3(aq)+\operatorname{Li}_3(bq)\right]
 +\frac13\ln^3q,\\
 G_{r,-1,-1}(X)&=-\frac12\left[\mathcal F(q)-\mathcal F(1)\right].
 \label{eq:comp-root-squared-classical}
\end{align}
Indeed, $\mathcal F'(q)=\ln^2[(1-q+q^2)/q]/q$.
The minus sign follows from $\omega_r=-\dd q/q$, and the
subtraction at $q=1$ enforces the original lower limit $X=0$.

The branches in this formula follow directly from the physical
interval. The arguments of $A$ and $B$ lie respectively in
$[-\pi/3,0]$ and $[0,\pi/3]$, so
$\log R=\log A-\log B$ throughout the interval.
The ratio $R$ follows the lower unit-circle arc from $1$ to
$e^{-2\pi\ii/3}$. Neither this arc nor the rotated arc $R/(a/b)$
crosses the polylogarithm cut $[1,\infty)$ in the open interval.
The other polylogarithm arguments are $A$, $B$, $aq$ and $bq$,
which also avoid that cut. The logarithms in $\mathcal J$ are
continuous on these paths. Thus principal values suffice, with no
numerically chosen branch corrections. Although individual terms
are complex, the derivative of $\mathcal F$ is real, so the endpoint
difference in \eqref{eq:comp-root-squared-classical} is real.
This reduction leaves five nonclassical words in the partial sum.

\paragraph{Classical forms for the five remaining words.}
The same primitive also integrates the mixed logarithms generated by
integration by parts. This gives a common classical representation for
all six words, rather than six unrelated integration formulas.
For $d\in\{a,b\}$ and $h\in\{1,-1\}$, define
\begin{align}
 B_d(q)&=1-dq,\qquad R_d(q)=q/B_d(q),\qquad
 T_{hd}(q)=(1-hq)/B_d(q),\\
 N_d(q)&=-\ln q\operatorname{Li}_2(dq)+\operatorname{Li}_3(dq),\\
 P_{hd}(q)&=\tfrac12\bigl[\mathcal J(1-hq,1)+\mathcal J(B_d,1)
             -\mathcal J(T_{hd},1)+\mathcal J(T_{hd},h/d)\bigr],\\
 H_{dc}(q)&=\tfrac12\bigl[\mathcal J(q,c)+\mathcal J(B_d,1-dc)
             -\mathcal J(R_d,c/(1-dc))+\mathcal J(R_d,-1/d)\bigr].
 \label{eq:comp-root-mixed-functions}
\end{align}
The definition of $N_d$ is also used for $d=1,-1$.
In $H_{dc}$ only $c=1,\ii,-\ii$ is needed.
To derive these expressions, apply the square-of-a-ratio identity
used above to each pair of logarithms. The transformed measures
include both the finite pole and the image of infinity. Direct
differentiation then gives
\begin{align}
 N_d'(q)&=\frac{\ln q\log(1-dq)}q,\\
 P_{hd}'(q)&=\frac{\log(1-hq)\log(1-dq)}q,\\
 H_{dc}'(q)&=\frac{\ln q\log(1-dq)}{q-c}.
 \label{eq:comp-root-mixed-derivatives}
\end{align}
These identities specify the primitives independently of any cross-section
formula from the literature.

Some individual terms have a singular expression at $q=1$, although
the required combinations are finite. Their boundary values must therefore
be evaluated before numerical subtraction. The necessary limits are
\begin{align}
 \mathcal J(0,\rho)&=0\quad(\rho\ne0),\qquad
 \mathcal J(1,1)=-2\zeta(3),\\
 P_{1d}(1)&=\tfrac12\mathcal J(1-d,1),\\
 H_{d1}(1)&=\tfrac12\bigl[-2\zeta(3)
       +\log^2(1-d)\log(-d)+4\zeta(2)\log(1-d)
       +\mathcal J(1/(1-d),-1/d)\bigr].
 \label{eq:comp-root-mixed-boundary}
\end{align}
For example, the two divergent logarithms in $H_{d1}$ occur in
$\mathcal J(B_d,1-d)-\mathcal J(R_d,1/(1-d))$.
As $q\to1^-$, their difference tends to $\log(-d)$.
Their coefficients both tend to $\log^2(1-d)$ because
$\log(1/(1-d))=-\log(1-d)$. The dilogarithms tend to
$\zeta(2)$, giving the displayed $4\zeta(2)\log(1-d)$ term.
This determines the boundary constant analytically.

Put $S(q)=N_a(q)+N_b(q)$ and define
\begin{align}
 U(q)&=-2\sum_{d=a,b}P_{1d}(q)+S(q)+2N_1(q)-\tfrac13\ln^3q,\\
 V(q)&=\ln q\operatorname{Li}_2(-X)
       +2\sum_{d=a,b}H_{d1}(q)-S(q)
       -2\mathcal J(q,1)+\tfrac13\ln^3q,\\
 W(q)&=2\sum_{d=a,b}P_{-1,d}(q)-S(q)-2N_{-1}(q)
       +\tfrac13\ln^3q+\ln4\,E(q),\\
 E(q)&=\operatorname{Li}_2(aq)+\operatorname{Li}_2(bq)
                 +\tfrac12\ln^2q,\\
 Z(q)&=-\ln q\,D(X)
       +\sum_{c=\ii,-\ii}\left[\sum_{d=a,b}H_{dc}(q)-\mathcal J(q,c)\right]
       -S(q)+\tfrac13\ln^3q,\\
 D(X)&=\operatorname{Li}_2(-1-X)+\ln(1+X)\ln(2+X)+\frac{\pi^2}{12}.
\end{align}
All occurrences of $X$ here mean $(1-q)^2/q$.
Write $\Delta f=f(q)-f(1)$, with the finite boundary values just derived.
The five reductions are
\begin{align}
 G_{r,0,-1}(X)&=\Delta V,\\
 G_{r,-1,0}(X)&=\Delta U-\Delta V,\\
 G_{0,r,-1}(X)&=\ln X\,\mathcal G(X)-\Delta U,\\
 G_{-4,r,-1}(X)&=\ln\!\left(\frac{X+4}{4}\right)\mathcal G(X)+\Delta W,\\
 G_{r,-2,-1}(X)&=\Delta Z.
 \label{eq:comp-root-five-classical}
\end{align}
Here $\mathcal G(X)=G_{r,-1}(X)$ is the classical dilogarithm
combination in \eqref{eq:comp-ir-short-function}.

For completeness, the integration steps fixing these formulas can be
checked without manipulating the full rational coefficients of the cross
section. Since $G_{0,-1}=-\operatorname{Li}_2(-X)$, integration by parts gives
\begin{equation}
 G_{r,0,-1}(X)=\ln q\operatorname{Li}_2(-X)
 +\int_1^q\ln t\ln(1+X(t))
       \left(\frac2{t-1}-\frac1t\right)\dd t.
\end{equation}
Inserting $\ln(1+X(t))=\log(1-at)+\log(1-bt)-\ln t$
and using \eqref{eq:comp-root-mixed-derivatives} gives $\Delta V$.
The shuffle identity gives
$G_{-1,0}+G_{0,-1}=\ln X\ln(1+X)$, whose integral with
$\omega_r$ is $\Delta U$. Multiplication of $G_{r,-1}$ by
$\ln X$ then gives the third line of
\eqref{eq:comp-root-five-classical}.
For the fourth line, integrate the outer $-4$ kernel by parts and use
\begin{equation}
 \ln\!\left(\frac{X(t)+4}{4}\right)
       =2\ln(1+t)-\ln t-\ln4.
\end{equation}
Finally, $D(X)=G_{-2,-1}(X)$ satisfies
$\dd D/\dd X=\ln(1+X)/(X+2)$. The factorization
$X(t)+2=(t-\ii)(t+\ii)/t$ gives
\begin{equation}
 \frac{\dd D(X(t))}{\dd t}
 =\ln(1+X(t))\left(\frac1{t-\ii}+\frac1{t+\ii}-\frac1t\right).
\end{equation}
Integrating $-D(X(t))\,\dd t/t$ by parts now gives $\Delta Z$.
Thus all constants follow from the original zero lower limit, and
no GPL remains among these six functions.

The complex arguments in the mixed primitives also have fixed physical
branches. For $d=a$, $B_d$ lies below the real axis, while
$q/B_d$ and $(1\pm q)/B_d$ lie above it. Their arguments stay
between $0$ and $\pi/3$, so each logarithm of a ratio is the
difference of the corresponding logarithms. The $d=b$ expressions
are their conjugate continuations. For $c=\pm\ii$, the arguments
$B_a/(1-ac)$ and $R_a(1-ac)/c$ can cross the real axis, but
only below $1$. This follows by multiplying numerator and denominator
by the conjugate denominator. For instance the possible crossing is
at $q=2-\sqrt3$, where $B_a/(1-a\ii)=2\sqrt3-3<1$ and
$R_a(1+a\ii)/(-\ii)=1-2\sqrt3/3<0$. Thus these paths do not cross the
polylogarithm cut. The sole real argument on that cut is $1+q$
in $\mathcal J(1+q,1)$. There the real boundary value is used:
\begin{equation}
 \mathcal J(z,1)=\ln^2z\ln(z-1)
 +2\ln z\operatorname{Re}\operatorname{Li}_2(z)
 -2\operatorname{Re}\operatorname{Li}_3(z),\qquad z>1.
\end{equation}
Its derivative is $\ln^2z/(z-1)$ and its limit at $z=1$ is
$-2\zeta(3)$. This fixes the continuation without a fitted imaginary
part. Equations \eqref{eq:comp-root-squared-classical} and
\eqref{eq:comp-root-five-classical} complete the classical reduction
of this virtual and UV partial sum. Further simplification of their
combined coefficients and arguments remains a separate step.

\paragraph{Combining the root terms in the cross section.}
The preceding conversion to classical functions does not by itself shorten
the cross section. Their coefficients must also be combined. Let
$C_{\mathrm{six}}(X)$ denote just the contribution of the six words listed
above to $\widehat C_{V+\mathrm{UV}}$. Their original rational coefficients
have the common factor $Y/[X^4(X+4)]$. After inserting their classical
forms, the coefficients of $\Delta U$ and $\Delta V$ are in the exact
ratio $-3:4$. Define $K(q)=4V(q)-3U(q)$ and use the polynomials in
\eqref{eq:comp-root-compact-polynomials}.
The six contributions then combine into
\begin{equation}
 \begin{aligned}
 C_{\mathrm{six}}(X)=\frac{Y}{X^4(X+4)}\bigg\{&
 -\frac32p_F\Delta\mathcal F+p_K\Delta K
 -6p_W\Delta W-8p_Z\Delta Z\\
 &+\left[3p_0\ln X-6p_W\ln\!\left(\frac{X+4}{4}\right)\right]
       \mathcal G(X)\bigg\}.
 \end{aligned}
 \label{eq:comp-root-combined-coefficients}
\end{equation}
Thus four boundary differences suffice for this contribution. In
particular, $U$ and $V$ need not be evaluated separately.
There is also a cancellation inside $K$: the $8\mathcal J(q,1)$
generated by $8\sum_d H_{d1}$ cancels its explicit
$-8\mathcal J(q,1)$ term. This follows by expanding
\eqref{eq:comp-root-mixed-functions}, before evaluating any special
function. Equation~\eqref{eq:comp-root-combined-coefficients} is an exact
coefficient collection of the six original terms in this diagram subset.
The other virtual, soft, and real terms must still be added using
the common-regulator assembly \eqref{eq:comp-assembly-common-regulator}.

\section{Combining the finite regulator conversions}
The three finite shifts in \eqref{eq:comp-simplified-real-remainder}
can be combined before evaluating the real contribution. Define
\begin{equation}
 M(X)=\Delta_{\mathrm{rat}}+\Delta_{\log}+\Delta_D,
 \qquad L=\ln(1+X),\quad J=\ln(1+X/2),\quad X>0.
\end{equation}
The shifts include the finite displacement of the double-pole
boundary in the Bose-restoring term. Combining their independently
integrated endpoint moments gives
\begin{equation}
 \begin{aligned}
 M(X)={}&-\frac{P_0}{3X}
 +\frac{2P_1}{X(1+X)}\operatorname{Li}_2(-X)
 -\frac{8(1+X)(X^2-4X-6)}{X^2}L^3\\
 &+\frac{16(1+X)(X^2-4X-6)}{X^2}L\operatorname{Li}_2(-X)
 -\frac{256(1+X)(X+2)^2}{X^3}J\\
 &+\frac{128(1+X)P_2}{3X^3}
 \left[\operatorname{Li}_2\!\left(-\frac{X^2}{4(1+X)}\right)-2JL\right]
 -\frac{P_3L^2+P_4L}{3X^3(1+X)}.
 \end{aligned}
 \label{eq:comp-matching-short}
\end{equation}
Here the polynomials are
\begin{align}
 P_0&=X^5-X^4-8X^3+196X^2+1104X+864,\\
 P_1&=X^5-6X^4-43X^3-32X^2+32X+32,\\
 P_2&=X^3-5X^2-28X-24,\\
 P_3&=3X^7-18X^6-233X^5-96X^4
                 +3008X^3+8240X^2+8000X+2688,\\
 P_4&=X^9+6X^8-18X^7-24X^6+363X^5
                 +1652X^4+412X^3-3984X^2-4704X-1536.
\end{align}
In particular, the explicit $\ln X$ cancels in this sum.
The reduction uses positive-real logarithm identities and
\begin{equation}
 \operatorname{Li}_2\!\left(\frac{X}{1+X}\right)
 =-\operatorname{Li}_2(-X)-\frac12\ln^2(1+X).
\end{equation}
The last identity follows by differentiating both sides and fixing
the zero value at $X=0$. Both sides are real for $X>0$.
The remaining dilogarithm arguments in
\eqref{eq:comp-matching-short} are negative, so their real branches
are fixed throughout the physical interval. Coefficient collection
then requires only rational identities in $X$. Thus the real remainder
can be written more compactly as
\begin{equation}
 \mathcal R_{\mathrm{rem}}(X)
 =F_{\mathrm{rat}}^z+F_W^z+F_{\log}^z+F_D^z
                  +2F_{UW}^{\eta}+M(X).
\end{equation}
The script \nolinkurl{compton_matching_simplify.py} constructs this
sum from the separate endpoint moments and verifies the rational
coefficients exactly. The overall factor $1/[8X(1+X)]$ is still
applied only in \eqref{eq:comp-simplified-master}.

\section{The finite soft contribution}
\label{sec:comp-ck-simplification}
The kernel was simplified before integration in
\eqref{eq:comp-ck-pole-removal}. Its real-letter endpoint expression is
\eqref{eq:comp-ck-simplified-endpoint}, with the integration constants
fixed by vanishing boundary terms. We now reduce that completed
expression to classical functions and then to the common energy basis.

\subsection{The finite soft coefficient in classical polylogarithms}
\label{sec:comp-ck-classical}
All weight-three terms in the preceding soft coefficient arise from
products of two ordinary logarithms.  This permits a classical
polylogarithmic endpoint expression for the complete coefficient,
without a GPL primitive operator.

For $z>0$ and real $\rho\ne0$, define
\begin{align}
 J(z,\rho)&=\ln^2z\ln|1-z/\rho|
 +2\ln z\operatorname{Re}\operatorname{Li}_2(z/\rho)
 -2\operatorname{Re}\operatorname{Li}_3(z/\rho),\\
 K(z,\rho)&=\ln z\ln|1-z/\rho|
 +\operatorname{Re}\operatorname{Li}_2(z/\rho).
\end{align}
The real parts denote the real boundary values on the standard
polylogarithm cut.  They enter because the original logarithms are real.
Their derivatives are listed below. At the degenerate values
$\rho=0,\infty$, the primitive is defined up to an additive constant,
which cancels in every normalized boundary difference. We use
\begin{align}
 \partial_z J(z,\rho)&=\frac{\ln^2z}{z-\rho},&
 J(z,0)&=\frac13\ln^3z,&J(z,\infty)&=0,\\
 \partial_z K(z,\rho)&=\frac{\ln z}{z-\rho},&
 K(z,0)&=\frac12\ln^2z.
\end{align}
The convention at infinity specifies a constant primitive with zero
derivative.  It will occur when a pole maps to infinity.

For a real letter $a$ outside the open physical interval, put
$A_a(z)=1-z/a$ and $\ell_a(z)=\ln A_a(z)$.  All factors used below
are positive for $0<z<Z$.  The normalized single-log primitive is
\begin{equation}
 \mathcal K_{a;c}(z)
 =K(A_a(z),A_a(c))-K(1,A_a(c)).
\end{equation}
For $a\ne b$ define
\begin{equation}
 \begin{aligned}
 \mathcal H_{a,b;c}(z)=\frac12\bigg[&
 J(A_a(z),A_a(c))+J(A_b(z),A_b(c))\\
 &-J\left(\frac{A_a(z)}{A_b(z)},
                 \frac{A_a(c)}{A_b(c)}\right)
 +J\left(\frac{A_a(z)}{A_b(z)},\frac ba\right)\bigg]
 -(z\to0).
 \end{aligned}
 \label{eq:comp-ck-classical-product}
\end{equation}
If $c=b$, the third term has its second argument at infinity and
vanishes.  If $a=b$, use
$\mathcal H_{a,a;c}(z)=J(A_a(z),A_a(c))-J(1,A_a(c))$.
These definitions obey
\begin{equation}
 \mathcal K'_{a;c}(z)=\frac{\ell_a(z)}{z-c},\qquad
 \mathcal H'_{a,b;c}(z)=\frac{\ell_a(z)\ell_b(z)}{z-c},\qquad
 \mathcal K_{a;c}(0)=\mathcal H_{a,b;c}(0)=0.
\end{equation}
To verify the product primitive, first use
$2\ln A\ln B=\ln^2A+\ln^2B-\ln^2(A/B)$.  Under
$v=A_a(z)/A_b(z)$, the rational differential transforms as
\begin{equation}
 \frac{\dd z}{z-c}
 =\frac{\dd v}{v-A_a(c)/A_b(c)}
  -\frac{\dd v}{v-b/a}.
\end{equation}
The second kernel is the image of infinity and is essential.
This identity derives both ratio terms in
\eqref{eq:comp-ck-classical-product}.  Positivity of $A_a,A_b$
justifies the logarithmic identity without additional multiples of
$\ii\pi$.  Subtraction at zero fixes every integration constant.

Here is a finite formula for including the rational coefficients.
For any rational $R(z)$ in the four soft kernels, write
\begin{equation}
 R=P_R'+\sum_c\frac{r_c(R)}{z-c},\qquad
 P_R=\sum_{n\geq0}\frac{p_nz^{n+1}}{n+1}
 -\sum_c\sum_{j\geq2}\frac{r_{c,j}}{(j-1)(z-c)^{j-1}}.
\end{equation}
The sums use its polynomial and partial-fraction coefficients.
They are finite algebraic derivatives of rational functions.
Let $\mathcal E[R](z)$ be the elementary primitive of $R$ normalized
to zero at the origin,
\begin{equation}
 \mathcal E[R](z)=P_R(z)-P_R(0)
 +\sum_cr_c(R)\ln\left|\frac{z-c}{-c}\right|.
\end{equation}
No pole here lies at the integration origin.  Define the following
explicit combinations of classical functions:
\begin{equation}
 \mathcal I_a[R]
 =P_R\ell_a+\sum_cr_c(R)\mathcal K_{a;c}
 -\mathcal E\!\left[\frac{P_R}{z-a}\right].
 \label{eq:comp-ck-classical-single}
\end{equation}
For the product, put $R_a=P_R/(z-a)$ and $R_b=P_R/(z-b)$.  Then
\begin{equation}
 \begin{aligned}
 \mathcal I_{a,b}[R]={}&P_R\ell_a\ell_b
 +\sum_cr_c(R)\mathcal H_{a,b;c}
 -P_{R_a}\ell_b-P_{R_b}\ell_a\\
 &-\sum_cr_c(R_a)\mathcal K_{b;c}
  -\sum_cr_c(R_b)\mathcal K_{a;c}\\
 &+\mathcal E\!\left[\frac{P_{R_a}}{z-b}
                         +\frac{P_{R_b}}{z-a}\right].
 \end{aligned}
 \label{eq:comp-ck-classical-double}
\end{equation}
These are nonrecursive formulas.  Direct differentiation gives
$\mathcal I_a'[R]=R\ell_a$ and
$\mathcal I_{a,b}'[R]=R\ell_a\ell_b$.
Their zero values at the origin follow from the normalized
$\mathcal E,\mathcal K,\mathcal H$ and $\ell_a(0)=0$.

Denote the weights of the three specified logarithm lists by
$b_a$ for $L_\beta$, $u_a$ for $L_1$ and $k_a$ for $L_k$.
The complete classical primitive is
\begin{equation}
 \begin{aligned}
 \Psi_{\mathrm{cl}}(z)={}&
 \sum_ab_a\mathcal I_a[\widetilde R_0]
 +\sum_{a,b}b_au_b\mathcal I_{a,b}[\widetilde R_1]\\
 &+\sum_{a,b}b_ak_b\mathcal I_{a,b}[\widetilde R_k]
 +\sum_ak_a\mathcal I_a[\widetilde R_f].
 \end{aligned}
\end{equation}
Each index ranges over its four explicitly given real letters.
Finally,
\begin{equation}
 \boxed{\mathcal C_K(s)=\lim_{z\to Z^-}\Psi_{\mathrm{cl}}(z)
 -\frac{\operatorname{artanh}\beta}{\beta}\mathcal C_0(s).}
 \label{eq:comp-ck-classical-complete}
\end{equation}
This expression uses only rational functions, real logarithms,
$\operatorname{Li}_2$ and $\operatorname{Li}_3$.  The endpoint limit
is taken after combining the terms.  It is the same finite limit
already established for the soft coefficient.

There is a further exact simplification before evaluation.  Each
logarithm list has weights summing to zero.  Hence the first two
$J$ terms in \eqref{eq:comp-ck-classical-product} cancel in both
weighted double sums, including their base-point constants.
Only the two ratio terms are needed there.  This cancellation
reduces the number of classical function evaluations without
changing the endpoint or its branch choice.

\subsection{Collection before the soft endpoint evaluation}
The soft coefficient can be shortened further by collecting the rational
integration-by-parts terms before expanding the logarithms into GPLs.
This also removes separate divergent boundary terms. Keep $X=s-1>0$,
$\beta=X/(X+2)$, and use the original angular variable $0<U<1$.
Define
\begin{align}
 k&=\sqrt{\frac{1+XU^2}{s}},\qquad d=1-\beta^2U^2,\\
 a&=\operatorname{artanh}(\beta U),\quad
 L=\operatorname{artanh}U,\quad f=\operatorname{artanh}(U/k).
\end{align}
The square root is positive. Since $k^2-U^2=(1-U^2)/s>0$, both
inverse hyperbolic tangents involving $U$ are real in the open interval.
Their derivatives are
\begin{equation}
 a'=\frac{\beta}{d},\qquad L'=\frac1{1-U^2},\qquad
 f'=\frac1{k(1-U^2)},\qquad k'=\frac{XU}{sk}.
\end{equation}
Primes here mean differentiation with respect to $U$.

\paragraph{Rational integration by parts.} Let $\Phi(U)$ be the soft
integrand after the pole removal in
\eqref{eq:comp-ck-pole-removal}, including its common factor
$4/(\beta E^2)$. Rational division and the derivatives above give
\begin{equation}
 \Phi(U)=\frac{\dd\mathcal B}{\dd U}+\sum_i T_i(U)f_i(U),
 \qquad (f_i)=(aL,af,a,L,f,1),
\end{equation}
where the rational boundary can be chosen as
\begin{equation}
 \mathcal B=Q_{aL}aL+Q_{af}af+Q_a a+Q_L L+Q_1.
\end{equation}
Its coefficients are
\begin{align}
 Q_{aL}&=\frac{2U(X+2)}{X^3s}
       \bigl(U^2X^3+X^3+24X^2+36X+16\bigr),\\
 Q_{af}&=\frac{128s^2Uk}{X^3(X+4)d},\\
 Q_a&=-\frac{2U^2(X+2)}s,\qquad
 Q_L=\frac{U^2(X+2)^2}{Xs},\\
 Q_1&=-\frac{U(X+2)^2}{Xs}.
\end{align}
For example, the root-dependent part follows directly from
\begin{equation}
 \frac{\dd}{\dd U}\frac{Uk}{d}
 =\frac{1+(2X+\beta^2)U^2}{skd^2}.
\end{equation}
Matching its constant and $U^2$ coefficients removes the double pole
from the coefficient of $af$. Substituting $f'$ then gives a rational
coefficient of $a$, because the factor $k$ in $Q_{af}$ cancels.
The other descendants follow by differentiating $aL$, $a$, and $L$.
Each step is a rational identity after replacing $k^2$ by
$(1+XU^2)/s$.

\paragraph{Vanishing logarithmic boundaries.} Replace each of
$Q_{aL},Q_{af},Q_a,Q_L$ by $Q_i(U)-Q_i(1)$ and add
$Q_i(1)f_i'$ to the remainder. This leaves the full integrand unchanged.
All these differences are $O(1-U)$. At the upper endpoint $a$ is finite,
while $L$ and $f$ grow logarithmically. Every logarithmic boundary
therefore vanishes. At $U=0$, all three logarithms vanish. The only
integrated boundary is
\begin{equation}
 Q_1(1)-Q_1(0)=-\frac{(X+2)^2}{Xs}.
\end{equation}
The cancellation of the apparent pole in the coefficient of $a$ uses
\begin{equation}
 \frac{1/k-1}{1-U^2}=\frac{X}{sk(1+k)}.
\end{equation}
This identity follows by multiplying by $1+k$ and using the displayed
expression for $k^2$. It also fixes the finite endpoint value without
subtracting two separately divergent numerical quantities.

Define the seven energy-dependent coefficients
\begin{align}
 A_L&=\frac{4(X+2)(X^3+12X^2+18X+8)}{X^3s},\\
 A_f&=\frac{32s(X+2)^2}{X^3(X+4)},\\
 P&=-\frac{16s(X^2-4X-6)}{X^2(X+2)},\\
 B&=\frac{16s(X^4+2X^3-24X^2-80X-48)}{X^3(X+2)(X+4)},\\
 C_L&=-\frac{4(X^3+14X^2+20X+8)}{X^2s},\\
 M&=-\frac{32s(X+2)}{X^2(X+4)},\qquad
 C_0=\frac{8(X+2)}{X^2}.
\end{align}
Here $A_L=Q_{aL}(1)$ and $A_f=Q_{af}(1)$.
The resulting remainder is the short expression
\begin{equation}
\begin{aligned}
 \widehat\Phi(U)={}&\frac1d\biggl[
 PaL+\frac{Baf}{k}
 +(A_L+A_f)\frac{1+\beta^2U}{1+U}a\\
 &\hspace{12mm}+(C_LU+A_L\beta)L
 +\left(\frac{MU}{k}+A_f\beta\right)f+C_0\beta\biggr]
 +\frac{A_fXa}{sk(1+k)}.
\end{aligned}
\label{eq:comp-ck-regular-root-kernel}
\end{equation}
None of its rational coefficients has a pole on the closed physical
interval. The logarithmic growth at $U=1$ is integrable. In particular,
the constant-function term integrates immediately to
\begin{equation}
 \int_0^1\frac{C_0\beta}{d}\,\dd U
 =C_0\operatorname{artanh}\beta
 =\frac{4(X+2)}{X^2}\ln s.
\end{equation}
Thus the already integrated soft coefficient admits the exact
reorganization
\begin{equation}
 \mathcal C_K(s)=-\frac{(X+2)^2}{Xs}
 +\int_0^1\widehat\Phi(U)\,\dd U
 -\frac{\operatorname{artanh}\beta}{\beta}\mathcal C_0(s).
\label{eq:comp-ck-regular-root-identity}
\end{equation}
Equation~\eqref{eq:comp-ck-regular-root-identity} is an identity for
collecting the existing analytic endpoint expression. Its remaining
terms are evaluated with the logarithmic primitives of the preceding
subsection. Under $U=(w-w^{-1})/(2\sqrt X)$, all their coefficients
are proper rational functions with simple poles, so only the single-pole
primitives are needed. The square root has not been expanded into
additional energy-dependent letters at the collection step.

\paragraph{Finite analytic endpoint sum.} The reorganized coefficient
can also be stated without a remaining integral. Put
\begin{equation}
 W=\sqrt s+\sqrt X,\qquad Z=W-1,\qquad
 r=\sqrt{1+X/\beta^2}-\frac{\sqrt X}{\beta}.
\end{equation}
Define three lists of pairs consisting of a letter and its weight,
\begin{align}
 \mathcal A&=\{(r,\tfrac12),(-r^{-1},\tfrac12),
                  (r^{-1},-\tfrac12),(-r,-\tfrac12)\},\\
 \mathcal L&=\{(W^{-1},\tfrac12),(-W,\tfrac12),
                  (W,-\tfrac12),(-W^{-1},-\tfrac12)\},\\
 \mathcal F&=\{(W^{-1},\tfrac12),(-W^{-1},\tfrac12),
                  (W,-\tfrac12),(-W,-\tfrac12)\}.
\end{align}
For example, $a(w)=\sum_{(b,h)\in\mathcal A}
h\ln[(w-b)/(1-b)]$. The other two lists give $L(w)$ and $f(w)$.
Use $U=(w-w^{-1})/(2\sqrt X)$ and $k=(w+w^{-1})/(2\sqrt s)$
in the five rational functions
\begin{align}
 R_{aL}&=U'\frac P d,\qquad R_{af}=U'\frac B{kd},\\
 R_a&=U'\left[\frac{(A_L+A_f)(1+\beta^2U)}{(1+U)d}
                      +\frac{A_fX}{sk(1+k)}\right],\\
 R_L&=U'\frac{C_LU+A_L\beta}{d},\qquad
 R_f=U'\frac{MU/k+A_f\beta}{d}.
\end{align}
Here $U'=\dd U/\dd w$. Each function is proper rational and has
only simple poles. If $\mathcal P_R$ is its pole set, its coefficients
are given by the finite algebraic limits
\begin{equation}
 R(w)=\sum_{c\in\mathcal P_R}\frac{\rho_c}{w-c},\qquad
 \rho_c=\lim_{w\to c}(w-c)R(w).
\end{equation}
For a list $\mathcal B$ and two lists $\mathcal B,\mathcal D$, define
the following finite sums of ordinary GPLs:
\begin{align}
 \mathcal S[R;\mathcal B]
 &=\sum_{c\in\mathcal P_R}\sum_{(b,h)\in\mathcal B}
       \rho_ch\,G(c-1,b-1;Z),\\
 \mathcal S[R;\mathcal B,\mathcal D]
 &=\sum_{c\in\mathcal P_R}
   \sum_{(b,h)\in\mathcal B}\sum_{(e,j)\in\mathcal D}
   \rho_chj\bigl[G(c-1,b-1,e-1;Z)
                    +G(c-1,e-1,b-1;Z)\bigr].
\end{align}
These follow by substituting the partial fractions and using the shuffle
identity for the product of the two logarithms. The soft coefficient is
therefore
\begin{equation}
\begin{aligned}
 \mathcal C_K(s)={}&-\frac{(X+2)^2}{Xs}
 +\frac{4(X+2)}{X^2}\ln s
 -\frac{\ln s}{2\beta}\mathcal C_0(s)\\
 &+\mathcal S[R_{aL};\mathcal A,\mathcal L]
  +\mathcal S[R_{af};\mathcal A,\mathcal F]\\
 &+\mathcal S[R_a;\mathcal A]
  +\mathcal S[R_L;\mathcal L]
  +\mathcal S[R_f;\mathcal F].
\end{aligned}
\label{eq:comp-ck-collected-finite-gpl}
\end{equation}
All the residues in this expression are algebraic functions specified
by the five short rational functions above. No endpoint extrapolation
or numerical integration is part of its definition.
The possible poles belong to
$\{0,\pm r,\pm r^{-1},W^{-1},-W,-W^{-1}\}$.
Since $0<r<W^{-1}<1<W<r^{-1}$, no outer letter lies on the
integration interval. The inner letter $W$ produces only the
integrable endpoint logarithms. Thus every GPL in this endpoint sum
is finite separately, and all logarithms follow the real physical
branch established above.

\subsection{The soft coefficient in the energy variable}
\label{sec:comp-ck-energy-formula}
The auxiliary endpoint letters of the preceding soft expression can be
removed entirely. The result uses the same energy variable $X$ and the
same root kernel as the virtual contribution. Throughout this subsection
$G(a_1,\ldots,a_n)$ means $G(a_1,\ldots,a_n;X)$, with
\begin{equation}
 \omega_a=\frac{\dd X}{X-a}\quad(a=0,-1,-2,-4),\qquad
 \omega_r=\frac{\dd X}{\sqrt{X(X+4)}}.
\end{equation}
The root is positive for $X>0$, and every function is continued from
$X=0$ along the positive real axis. Set $s=1+X$ and
$Y=\sqrt{X/(X+4)}$ as above. Define the combinations
\begin{align}
 \mathcal D&=G(0,-1)+G(-2,-1),\\
 \mathcal V&=G(-1,-1,-1)-G(-1,-2,-1)-G(-1,0,-1),\\
 \mathcal R&=3G(0,r,-1)-3G(-1,r,-1)-3G(-4,r,-1)
              +4G(r,-2,-1)+4G(r,0,-1).
\end{align}
Then the complete finite soft coefficient is
\begin{equation}
 \boxed{\mathcal C_K(1+X)=c_0+c_1G(-1)+c_{11}G(-1,-1)
           +c_D\mathcal D+c_V\mathcal V+c_rG(r,-1)+c_R\mathcal R.}
 \label{eq:comp-ck-energy-compact}
\end{equation}
The seven coefficients are
\begin{align}
 c_0&=-\frac{(X+2)^2}{Xs},\\
 c_1&=-\frac{(X+2)(X^3+14X^2+28X+16)}{X^3s},\\
 c_{11}&=-\frac{(X+2)(7X^4+32X^3+24X^2-16X-16)}{X^4s},\\
 c_D&=\frac{2(X+2)(X^4+24X^3+98X^2+120X+48)}{X^3s(X+4)},\\
 c_V&=\frac{4s(X^2-4X-6)}{X^3},\\
 c_r&=-\frac{24Ys(X+2)^3}{X^4(X+4)},\\
 c_R&=\frac{2Ys(X^4+2X^3-24X^2-80X-48)}{X^4(X+4)}.
\end{align}
This formula contains fourteen GPL words when the three combinations
are expanded, including the constant term. All coefficients are
rational in $X$ and $Y$. In particular, the auxiliary root $\sqrt s$
and all separate $\ln2$ terms have disappeared. In the total coefficient
\eqref{eq:comp-simplified-master}, this expression enters as
$-\mathcal C_K(s)/s$.

\paragraph{Reduction of the lower-weight integrals.} Use the functions
$a,L,f,k,d$ in the preceding subsection and define
\begin{align}
 I_1&=\int_0^1\frac{aL}{d}\,\dd U,&
 I_2&=\int_0^1\frac{af}{kd}\,\dd U,\\
 I_3&=\int_0^1\frac{a(1+\beta^2U)}{(1+U)d}\,\dd U,&
 I_4&=\int_0^1\frac{a}{k(1+k)}\,\dd U,\\
 I_5&=\int_0^1\frac{UL}{d}\,\dd U,&
 I_6&=\int_0^1\frac{L}{d}\,\dd U,\\
 I_7&=\int_0^1\frac{Uf}{kd}\,\dd U,&
 I_8&=\int_0^1\frac{f}{d}\,\dd U.
\end{align}
These integrals organize the derivation of the closed expression, rather
than define functions that remain to be integrated. Put
$a_*=a(1)=\tfrac12\ln s$. Since
$(f-L)'=X/[sk(1+k)]$, integration by parts gives
\begin{equation}
 \frac Xs I_4+\beta I_8=a_*^2+\beta I_6.
 \label{eq:comp-ck-lower-root-cancellation}
\end{equation}
At $U=0$ the boundary vanishes. At $U=1$ the finite difference
$f-L$ tends to $\tfrac12\ln s=a_*$, because
$1-U/k=(1-U)/s+O((1-U)^2)$. This determines the finite boundary
constant in \eqref{eq:comp-ck-lower-root-cancellation}.

The rational-letter integrals are
\begin{align}
 \beta I_1&=-\tfrac14\mathcal V,\\
 I_3&=\tfrac12G(0,-1)-\tfrac12\ln2\,G(-1)
       -\tfrac12G(-1,-2)-\tfrac14G(-1,-1)+\tfrac12G(-2,-1),\\
 \beta^2 I_5&=\tfrac14G(-1,-1),\\
 \beta I_6&=\tfrac12\ln2\,G(-1)-\tfrac14G(-1,-1)
                    +\tfrac12G(-1,-2).
\end{align}
To obtain these expressions, factor $d$ into its two linear factors
and write $a$ and $L$ as differences of logarithms. The resulting GPLs
have endpoint one and letters $\pm1/\beta,\pm1$. Their total
differentials with respect to $X$ contain only the logarithmic
differentials of $X$, $X+1$, and $X+2$. Integration of those
differentials gives the displayed energy-variable words. Every word
containing a moving letter vanishes as $X\to0$, which fixes its
integration constant. An arbitrary tangential constant for the inner
letter one cancels before the result is collected. Thus no value is
assigned to a divergent logarithm in obtaining these identities.

\paragraph{Energy equations for the root integrals.} The two remaining
integrals $I_2$ and $I_7$ can be integrated by differentiating with
respect to $X$. At fixed $U$,
\begin{equation}
 \partial_X a=\frac{2U}{(X+2)^2d},\qquad
 \partial_X f=\frac{U}{2sk},\qquad
 \partial_X k^2=\frac{U^2-1}{s^2}.
\end{equation}
For any compact positive interval of $X$, these derivatives have at
most integrable logarithmic endpoint behavior. Differentiation under
the integral is therefore allowed. Introduce
\begin{align}
 Q_7&=p_7\frac{k}{d}+h_7\frac1k,&
 p_7&=\frac{2(X+2)}{X^2(X+4)},&
 h_7&=-\frac{(X+2)^2}{2X^2s(X+4)},\\
 Q_2&=U\left(p_2\frac{k}{d}+h_2\frac1k\right),&
 p_2&=\frac2{(X+2)(X+4)},&
 h_2&=\frac{(X+2)^2}{2Xs(X+4)}.
\end{align}
Direct differentiation gives the two rational identities
\begin{align}
 \partial_X\frac{U}{kd}
 &=\partial_UQ_7-\frac{2(X+6)}{X(X+2)(X+4)}\frac{U}{kd},\\
 \partial_X\frac1{kd}
 &=\partial_UQ_2-\frac4{X(X+2)(X+4)}\frac1{kd}.
\end{align}
They are verified by substituting $k^2=(1+XU^2)/s$ and matching
rational coefficients. For the first equation, integrate the total
derivative against $f$ using $Q_7-Q_7(1)$, so its boundary vanishes.
The remaining rational terms and the integral of $f'-L'$ give
\begin{equation}
 I_7'=-\frac{2(X+6)}{X(X+2)(X+4)}I_7
       +\frac{3(X+2)^2}{2X^2(X+4)}a_*.
\end{equation}
The integrating factor is $X\sqrt{X(X+4)}/(X+2)^2$, hence
\begin{equation}
 \frac{\dd}{\dd X}
 \left[\frac{X\sqrt{X(X+4)}}{(X+2)^2}I_7\right]
 =\frac34\frac{\ln(1+X)}{\sqrt{X(X+4)}}.
\end{equation}
At $X=0$, $I_7\to\int_0^1U\operatorname{artanh}U\,\dd U=1/2$,
so the product on the left vanishes. Therefore
\begin{equation}
 I_7=\frac{3(X+2)^2}{4X\sqrt{X(X+4)}}G(r,-1).
 \label{eq:comp-ck-I7-energy}
\end{equation}

For the second equation the endpoint cancellation is implemented by
subtracting $Q_2(1)a_*$ from $Q_2a$ before multiplying by $f'$.
Define the two lower-weight integrals
\begin{equation}
 J_{Ua}=\int_0^1\frac{Ua}{d}\,\dd U,\qquad
 J=\int_0^1\frac{a_*-Ua}{1-U^2}\,\dd U.
\end{equation}
The subtracted numerator in $J$ vanishes at the upper endpoint.
The differential equation becomes
\begin{equation}
\begin{aligned}
 I_2'={}&-\frac4{X(X+2)(X+4)}I_2
 -\frac{X-2}{2s(X+4)}I_7\\
 &+\frac X{(X+2)(X+4)}J_{Ua}
 +\frac{X+2}{X(X+4)}(J+a_*^2).
\end{aligned}
\end{equation}
The rational-letter reductions give
\begin{align}
 \beta^2J_{Ua}&=-\tfrac14G(-1,-1)+\tfrac12G(-2,-1),\\
 J&=\tfrac12G(0,-1)-\tfrac14G(-1,-1),\\
 \beta^2J_{Ua}+J+a_*^2&=\tfrac12\mathcal D.
\end{align}
Using \eqref{eq:comp-ck-I7-energy} and the integrating factor
$\gamma=\sqrt{X(X+4)}/(X+2)$ now yields
\begin{equation}
 \frac{\dd(\gamma I_2)}{\dd X}
 =\frac38\left(\frac1X-\frac1{X+1}-\frac1{X+4}\right)G(r,-1)
       +\frac{\mathcal D}{2\sqrt{X(X+4)}}.
\end{equation}
Each term is a defining GPL differential. Since $I_2=O(X)$ at small
$X$, the constant at $X=0$ is zero, and
\begin{equation}
 \gamma I_2=\frac18\mathcal R.
 \label{eq:comp-ck-I2-energy}
\end{equation}
Substituting these evaluated integrals into
\eqref{eq:comp-ck-regular-root-identity} proves
\eqref{eq:comp-ck-energy-compact}. The two terms proportional to
$A_f$ in \eqref{eq:comp-ck-lower-root-cancellation} cancel the
remaining $G(-1,-2)$ and $\ln2$ terms exactly.

\section{A shorter finite UW expression}
\label{sec:comp-UW-simplified}
The finite $UW$ contribution already has the integrated endpoint form
\eqref{eq:comp-real-UW-final-primitive}. Its thirteen inner monomials
contain repeated weight-two functions. Combining them before the
rational primitive reduces the endpoint expression. Throughout this
section $H=s-1>0$, the electron mass is one, and $0<t<1$.
Define
\begin{equation}
 P=1+Ht-Ht^2,\qquad \ell=\ln(1+Ht),\qquad
 p=\ln P,\qquad b=\ln(1+Ht/2).
\end{equation}
All three logarithms are real. The two weight-two combinations needed
below are
\begin{align}
 J_2(H,t)&=\operatorname{Li}_2(-1-Ht)+\frac{\pi^2}{12}
       +\ell\ln(2+Ht),\\
 J(H,t)&=\operatorname{Li}_2(t/P)-\operatorname{Li}_2(t)
       -p\ln(1-t)+\frac{p^2}{2}.
 \label{eq:comp-UW-J-short}
\end{align}
The dilogarithm arguments in $J$ obey $0<t/P\leq t<1$.
The argument in $J_2$ is negative. These expressions therefore fix the
physical real branches directly.

\paragraph{Exact coefficient combinations.} Let $r_f(H,t)$ denote the
coefficient of the atom monomial $f$ in the thirteen-term inner result.
Its rational coefficients obey
\begin{equation}
 r_{d_3}=-r_{\ell\ln t}=r_{\ell p},\qquad
 r_{d_2}=r_{\ell\ln(2+Ht)},\qquad
 [\ln2]r_1=-r_{\ln(2+Ht)}.
 \label{eq:comp-UW-coefficient-combinations}
\end{equation}
The last notation means the coefficient of $\ln2$ in $r_1$.
These are polynomial identities after multiplication by their common
rational denominators, for arbitrary $H$ and $t$. Here
$d_2=\operatorname{Li}_2(-1-Ht)+\pi^2/12$ and $d_3$ is the two-dilogarithm
difference defined in the inner integration. Thus its finite kernel is
\begin{equation}
 \Phi_{UW}=A_0+A_1\ell+A_2\ell^2+A_p p+A_b b
             +A_d\operatorname{Li}_2(-Ht)+A_{J_2}J_2+A_JJ,
 \label{eq:comp-UW-eight-kernel}
\end{equation}
with the following factored sums of the original rational coefficients:
\begin{align}
 A_0&=\left.r_1\right|_{\ln2=0}+r_{\ln H}\ln H,
 &A_1&=r_\ell+r_{\ell\ln H}\ln H,\\
 A_2&=r_{\ell^2},& A_p&=r_p,\\
 A_b&=r_{\ln(2+Ht)},& A_d&=r_{d_1},\\
 A_{J_2}&=r_{d_2},& A_J&=r_{d_3}.
\end{align}
Here $d_1=\operatorname{Li}_2(-Ht)$. Keeping these sums factored avoids
expanding large common numerators. The explicit eight coefficients are
stored with the calculation in
\texttt{\detokenize{uw_simplified_kernel.json}}.

\paragraph{Dilogarithm reduction.} To obtain $J$ without guessing a
branch, first apply the negative-real-axis identity
$\operatorname{Li}_2(-z)+\operatorname{Li}_2(-1/z)
=-\pi^2/6-\ln^2z/2$, with $z>0$, to $d_3$.
It gives the analytic combination $D=d_3-\ell\ln t$ derived at the lower
endpoint. Next use
\begin{equation}
 \operatorname{Li}_2(-z)
 =-\operatorname{Li}_2\!\left(\frac{z}{1+z}\right)
  -\frac12\ln^2(1+z),\qquad z>0.
\end{equation}
The two positive arguments are
$t/[(1-t)(1+Ht)]$ and $t/(1-t)$.
Their transformed arguments are $t/P$ and $t$. The logarithms combine
exactly to give $J=D+\ell p$, hence \eqref{eq:comp-UW-J-short}.
In particular no explicit $\ln t$ survives. Both $J(0)=0$ and $J_2(0)=0$.
At the upper endpoint their analyticity follows from the original
$d_3-\ell\ln t+\ell p$ and $d_2+\ell\ln(2+Ht)$ representations.
This observation also fixes the value at $t=1$, where separate terms
of \eqref{eq:comp-UW-J-short} have logarithmic derivatives.

\paragraph{The integrated result.} Introduce the letters
\begin{equation}
 a=-\frac1H,\qquad c=-\frac2H,\qquad
 r_\pm=\frac{1\pm\sqrt{1+4/H}}2.
\end{equation}
In the same GPL convention as the rational primitive,
\begin{align}
 J_2(H,t)&=G(c,a;t),\\
 J(H,t)&=\sum_{r=r_\pm}
       \bigl[G(0,r;t)-G(1,r;t)+G(r,a;t)\bigr]-G(0,a;t).
\end{align}
For example $J_2'=H\ell/(2+Ht)$ and $J_2(0)=0$ prove its single-word
form. The formula for $J$ follows from the shuffle product
$G(a;t)G(r;t)=G(a,r;t)+G(r,a;t)$ and the expression for $D$.
The complete finite integral is therefore
\begin{equation}
\begin{aligned}
 \mathcal F_{UW}(H)=\lim_{Z\to1^-}\bigg\{&
 \mathfrak T[A_0;\varnothing](Z)+\mathfrak T[A_1;a](Z)
 +2\mathfrak T[A_2;a,a](Z)\\
 &+\mathfrak T[A_b;c](Z)+\mathfrak T[A_{J_2};c,a](Z)
 -\mathfrak T[A_d+A_J;0,a](Z)\\
 &+\sum_{r=r_\pm}\Bigl(
 \mathfrak T[A_p;r](Z)+\mathfrak T[A_J;0,r](Z)\\
 &\hspace{26mm}-\mathfrak T[A_J;1,r](Z)
 +\mathfrak T[A_J;r,a](Z)\Bigr)\bigg\}.
\end{aligned}
\label{eq:comp-UW-short-primitive}
\end{equation}
The operator $\mathfrak T$ is the finite rational-to-GPL primitive
already evaluated by the terminating rules above, including its
zero finite part at the lower endpoint. Thus this is a shorter
representation of the completed integral, with fourteen primitive
contributions. It is not a replacement by an unevaluated quadrature.
The terms with first letter one must be summed before taking the
upper limit. The total cross section uses this sector with the same
factor $1/(8sH)$ and the second Bose-related contribution as before.

\paragraph{Removal of apparent endpoint poles.} The certificates in
Sections~\ref{sec:comp-UW-upper-certificate} and
\ref{sec:comp-UW-lower-certificate} also give a form in which each
summand is finite at both endpoints. For any one of the eight functions
$f_i$ in \eqref{eq:comp-UW-eight-kernel}, decompose its coefficient as
\begin{equation}
 A_i(t)=A_i^{\mathrm{reg}}(t)
       +\sum_{e=0,1}\sum_{k=1}^{m_{ie}}\frac{c_{iek}}{(t-e)^k}.
\end{equation}
The regular coefficient has no pole at zero or one. Since every $f_i$
is analytic there, define its finite Taylor subtraction
\begin{equation}
 f_i^{[e,k]}(t)=\frac{f_i(t)-\displaystyle\sum_{n=0}^{k-1}
                  f_i^{(n)}(e)(t-e)^n/n!}{(t-e)^k}.
\end{equation}
Then the same finite kernel satisfies the exact identity
\begin{equation}
 \Phi_{UW}(t)=\sum_i A_i^{\mathrm{reg}}(t)f_i(t)
          +\sum_{i,e,k}c_{iek}f_i^{[e,k]}(t).
\end{equation}
Indeed, the omitted Taylor pieces are exactly the principal parts
of the full kernel at zero and one, which the two symbolic certificates
show to vanish. These divided differences are explicit logarithms and
dilogarithms minus finite polynomials. They introduce no new integral.
The factored form \eqref{eq:comp-UW-short-primitive} is shorter for
presenting the integrated result.

\section{Reciprocal symmetry of the real-emission integral}
After the first energy integration, write $s=1+H$, with $H>0$,
and $z=s+1-u-s/u$. The two algebraic boundaries of the first
integral are $x_-=s-u$ and $x_+=s-s/u$. Under $u\mapsto s/u$,
$z$ is unchanged and these boundaries exchange places. The
oriented boundary difference therefore changes sign. Its
Jacobian $j=1-s/u^2$ also changes sign, so the complete last
integrand obeys
\begin{equation}
 \frac{s}{u^2}f(s/u)=f(u).
\end{equation}
This identity holds separately for the regular rational,
energy-logarithm, $W$, and Bose-restoring sectors. It can also be
checked coefficient by coefficient: $L_1,L_2,L_3$ change sign,
whereas $L_1L_2$ is invariant under the reciprocal transformation.

The physical upper branch has $\sqrt{s}<u<s$. Its reflected
partner is $1<u<\sqrt{s}$. A common cut $z>\epsilon$ selects
reciprocal endpoints $u_-(\epsilon)u_+(\epsilon)=s$. Consequently
\begin{equation}
 \int_{\sqrt{s}}^{u_+(\epsilon)}f(u)\,\dd u
 =\frac12\int_{u_-(\epsilon)}^{u_+(\epsilon)}f(u)\,\dd u.
\end{equation}
Thus the square-root lower boundary can be eliminated exactly,
without changing the infrared subtraction or making an asymptotic
approximation.

Let $A$ be the coefficient of $1/(s-u)$ at the upper endpoint.
Reciprocal symmetry makes it also the coefficient of $1/(u-1)$
at the lower endpoint. Define
\begin{equation}
 \operatorname{FP}_{[1,s]}\int f
 =\lim_{\delta\to0^+}
 \left[\int_{1+\delta}^{s-\delta}f(u)\,\dd u+2A\ln\delta\right].
\end{equation}
Since $z=H(u-1)+O((u-1)^2)$ at the lower boundary and
$z=H(s-u)/s+O((s-u)^2)$ at the upper boundary, the finite part
in the original $z$ convention is
\begin{equation}
 F_z=\frac12\operatorname{FP}_{[1,s]}\int f
                  +A\ln\frac{H}{\sqrt{s}}.
\end{equation}
The explicit logarithm is essential. The replacement of the lower
boundary does not permit changing the finite subtraction constant.

The Bose-restoring term also has a double endpoint pole. If
\begin{equation}
 f(s-b)=D/b^2+A/b+O(1),
\end{equation}
reciprocal symmetry gives $f(1+a)=D/(sa^2)+A/a+O(1)$.
The finite-part definition then also subtracts
$D(1+1/s)/\delta$. Expanding both cut boundaries to second order
gives the corrected general formula
\begin{equation}
 F_z=\frac12\operatorname{FP}_{[1,s]}\int f
       +A\ln\frac{H}{\sqrt{s}}-\frac{D(s+1)}{2sH}.
\end{equation}
The previous expression is the case $D=0$. The extra constant
comes from the double pole multiplied by the second-order change
of the endpoint, and must be kept for the Bose term.

\paragraph{Eliminating the apparent threshold at $s=9$.}
The regular rational primitive contains ten words whose first letters are
\begin{equation}
 q_\pm=\frac{s-3\pm\sqrt{(s-1)(s-9)}}2.
\end{equation}
They can be replaced by a single real dilogarithm and GPLs whose letters
already occur in the other terms. This replacement follows from the
coefficients of our primitive. It does not use a total-cross-section formula.
Throughout this paragraph the GPL basepoint is $r=\sqrt{s}$, as in the
original primitive. Define
\begin{align}
 A(u)&=\frac{(u-1)(u+s)^2}{u(u+1)^2(s-u)},
 &Q(u)&=u^2-(s-3)u+s,\\
 (k_0,k_1,k_s,k_{-s},k_{-1})&=(-1,1,-1,2,-2),
 &C(s)&=\frac{128s(s+5)}{3(s-1)}.
\end{align}
The ten coefficients are precisely $C(s)k_b$ for each of the two first
letters $q_\pm$ and each $b\in\{0,1,s,-s,-1\}$. Since $A(r)=1$,
\begin{equation}
 L_3(u)=\ln A(u)=\sum_b k_bG(b;u).
\end{equation}
Factorization of $1-A$ gives the rational differential identity
\begin{equation}
 \frac{Q'(u)}{Q(u)}
 =\frac{\dd}{\dd u}\ln(1-A(u))
 -\frac1{u-r}-\frac1{u+r}
 +\frac1u+\frac1{u-s}+\frac2{u+1}.
\end{equation}
Here the logarithmic derivative denotes the rational quotient
$-A'(u)/(1-A(u))$ wherever its separate terms have removable poles.
Multiplying by $L_3$ and using
$\dd\operatorname{Li}_2(1-A)=-\ln A\,\dd\ln(1-A)$ yields
\begin{align}
 \sum_{\epsilon=\pm}\sum_b k_bG(q_\epsilon,b;u)
 ={}&-\operatorname{Li}_2(1-A(u))\\
 &+\sum_b k_b\bigl[-G(r,b;u)-G(-r,b;u)
       +G(0,b;u)\bigr]\\
 &+\sum_b k_b\bigl[G(s,b;u)+2G(-1,b;u)\bigr].
 \label{eq:comp-real-qpair-reduction}
\end{align}
Both sides vanish at $u=r$, which fixes the integration constant.
Although $r$ is a first letter on the right, the summed second-letter
function $L_3$ vanishes there, so the corresponding lower endpoint is
regular. For $r<u<s$, $A(u)>0$. Thus $1-A(u)<1$, and the dilogarithm
in this expression lies on its real branch. At a zero of $Q$, $A=1$
and this dilogarithm is regular. The formula also holds by continuity
at $s=9$, without choosing a side of an artificial threshold.

Replacing the ten words by \eqref{eq:comp-real-qpair-reduction} removes
$q_\pm$ from the primitive. It reintroduces the letters $\pm r$, so it
is a trade between two representations rather than a simultaneous
elimination of every auxiliary root. The accompanying coefficient
record contains the full replaced primitive. Its ten input coefficients,
rational differential identity, and lower-endpoint constant are checked
exactly by \texttt{compton\_real\_equivalence\_qpair.py}.

\paragraph{Combining the rational and $W$ primitives.}
The two primitives should be added before selecting a function basis.
Their original lists contain 48 and 19 nonempty GPL words. Collecting
identical words and canceling their exact rational coefficients leaves
43 words. In particular, the six Gram words discussed above disappear.
The coefficients in this combined list can be reduced further by shuffle
identities. For weight two,
\begin{equation}
 G(a,b)+G(b,a)=G(a)G(b),\qquad G(a,a)=\frac12G(a)^2.
\end{equation}
Applying these identities only to pairs present in the combined list leaves
40 GPL words and three elementary logarithm products.

There is a further reduction in the weight-three part. Define
\begin{equation}
 K=\frac{32s(2s+1)}{(s-1)^2},\qquad g=G(0;u)=\ln(u/\sqrt{s}).
\end{equation}
The entire weight-three contribution of the combined primitive is
\begin{align}
 K\bigl[&G(0,0,s)+G(0,s,0)-G(0,0,1)-G(0,1,0)
          -2G(0,0,0)\bigr]\\
 &=K\bigl[g\{G(0,s)-G(0,1)\}
          -G(0,0,s)+G(0,0,1)-g^3/3\bigr].
\end{align}
This follows from
$G(0)G(0,b)=2G(0,0,b)+G(0,b,0)$ for $b=1,s$.
Consequently the combined representation requires 37 nonempty GPL words,
three logarithm products, two logarithm--weight-two products, and one
logarithm cube. These are representation counts, not a claim of functional
independence of all remaining GPLs. The rational term is included throughout.

By contrast, substituting \eqref{eq:comp-real-qpair-reduction} before the
same combination gives 50 GPL words. Weight-two shuffle reduction leaves
42 words, eight logarithm products, and the additional dilogarithm.
The root-pair identity is therefore useful for displaying regularity at
$s=9$, but the original root-pair representation gives the shorter combined
formula. The full coefficient dictionary for the 37-word form and exact
reconstruction checks are recorded by
\texttt{compton\_real\_equivalence\_combine.py}.

\paragraph{Classical trilogarithms for the weight-three term.}
For positive real $z$, define $\mathcal L_n(z)$ here as
$\operatorname{Re}\operatorname{Li}_n(z+\ii0)$, for $n=2,3$.
The two boundary values are complex conjugates, so their real parts agree.
The differential relations are
\begin{equation}
 \frac{\dd\mathcal L_2(z)}{\dd z}=-\frac{\ln|1-z|}{z},
 \qquad
 \frac{\dd\mathcal L_3(z)}{\dd z}=\frac{\mathcal L_2(z)}{z}.
\end{equation}
They determine the continuation needed on each of the intervals used below.
For $b=1,s$ and $r<u<s$, direct differentiation and the zero value at $u=r$
give
\begin{align}
 G(0,b;u)&=-\mathcal L_2(u/b)+\mathcal L_2(r/b)
                      -\ln|1-r/b|\,g,\\
 G(0,0,b;u)&=-\mathcal L_3(u/b)+\mathcal L_3(r/b)
       +\mathcal L_2(r/b)g-\tfrac12\ln|1-r/b|\,g^2.
\end{align}
No singularity at $u=b$ is crossed along either open integration interval.
The complete weight-three contribution therefore reduces to
\begin{align}
 \mathcal F^{(3)}_{R_0+W}(u)=K\bigl[&
 \mathcal L_3(u/s)-\mathcal L_3(u)
 -\mathcal L_3(1/r)+\mathcal L_3(r)\\
 &+g\{\mathcal L_2(u)-\mathcal L_2(u/s)\}
 +\tfrac12g^2\ln r-\tfrac13g^3\bigr].
 \label{eq:comp-real-classical-weight-three}
\end{align}
The logarithmic coefficient uses
$|1-r|/|1-r/s|=r$. Thus the five original weight-three words require
only classical trilogarithms, dilogarithms, and logarithms. The lower
endpoint still vanishes exactly.

The remaining weight-two GPLs admit a uniform classical representation.
For distinct letters $a,b$ write
\begin{equation}
 z_a(u;b)=\frac{a-u}{a-b},\qquad z_a(r;b)=\frac{a-r}{a-b}.
\end{equation}
Then
\begin{align}
 G(a,b;u)={}&-\operatorname{Li}_2(z_a(u;b))
              +\operatorname{Li}_2(z_a(r;b))\\
 &-\operatorname{Log}(1-z_a(r;b))G(a;u).
 \label{eq:comp-real-classical-weight-two}
\end{align}
The logarithm and dilogarithm on the right must be continued together
from $u=r$ along the same path as the original GPL. More explicitly, their
logarithms satisfy
\begin{equation}
 \operatorname{Log}(1-z_a(u;b))
 -\operatorname{Log}(1-z_a(r;b))=G(b;u).
\end{equation}
This condition prevents independent principal-branch substitutions at the
two endpoints. Differentiation of \eqref{eq:comp-real-classical-weight-two}
then gives $G(b;u)/(u-a)$, and its initial value is zero. It therefore
proves the identity, including the continuation around removable letters.
Conjugate letters are combined before taking the real physical value.
Together with \eqref{eq:comp-real-classical-weight-three}, this eliminates
all GPL notation from the combined primitive if a classical-function
representation is preferred. Expanding every such identity separately
need not shorten the displayed expression, so conjugate pairs and common
coefficients should remain grouped.

At the upper endpoint the weight-three part simplifies further.
Put $\ell=\tfrac12\ln s$. The real inversion identities for $z>1$ are
\begin{align}
 \mathcal L_2(z)+\operatorname{Li}_2(z^{-1})
 &=\frac{\pi^2}{3}-\frac12\ln^2 z,\\
 \mathcal L_3(z)-\operatorname{Li}_3(z^{-1})
 &=\frac{\pi^2}{3}\ln z-\frac16\ln^3z.
\end{align}
Their derivatives agree on $z>1$, and their constants follow from
$\operatorname{Li}_2(1)=\pi^2/6$ and
$\operatorname{Li}_3(1)=\zeta(3)$. Applying them at $z=s,r$ gives
\begin{equation}
 \mathcal F^{(3)}_{R_0+W}(s)
 =K\left[\zeta(3)-\operatorname{Li}_3(s^{-1})
 -\ell\operatorname{Li}_2(s^{-1})
 -\frac{\pi^2\ell}{6}-\frac{2\ell^3}{3}\right].
 \label{eq:comp-real-classical-weight-three-endpoint}
\end{equation}
The lower endpoint is zero. The resulting expression depends only on
$s$ and contains no auxiliary square root or continued polylogarithm.

\paragraph{Finite energy endpoints before further combination.}
For the combined rational and $W$ primitive, cancellation of common
factors makes every rational coefficient regular at $u=s$. Thus its
finite upper endpoint needs only the coefficients evaluated at $s$.
No derivative of a coefficient multiplying a power-divergent GPL is
required. This statement follows by exact cancellation of their rational
denominators, whose endpoint factors are powers of $s$, $s-1$, and $s+3$.

For a one-letter word define its constant endpoint term by
\begin{equation}
 h_a=\begin{cases}
 \operatorname{Log}[(s-a)/(r-a)],&a\ne s,\\
 -\ln(s-r),&a=s.
 \end{cases}
\end{equation}
The second line removes the term $\ln(s-u)$ from $G(s;u)$.
For a weight-two word write $d_{ab}$ for the corresponding constant term.
For $a\ne s$, it is the boundary value of
\eqref{eq:comp-real-classical-weight-two}. For $a=s$ it is explicitly
\begin{equation}
 d_{sb}=\operatorname{Li}_2\left(\frac{s-r}{s-b}\right)
 +\operatorname{Log}\left(\frac{r-b}{s-b}\right)\ln(s-r).
\end{equation}
Indeed, the omitted divergent term is
$-\operatorname{Log}[(r-b)/(s-b)]\ln(s-u)$.
All logarithms use the continuation specified above. The only
logarithm products in the reduced primitive are $G(0)^2$, $G(0)G(1)$,
and $G(0)G(s)$, so their constant terms are simply the corresponding
products of the $h_a$. There is no product of two divergent one-letter
functions in this representation.

Replacing each word in the coefficient table by these constant endpoint
terms, and replacing its entire weight-three block by
\eqref{eq:comp-real-classical-weight-three-endpoint}, gives a function
$E_{R_0+W}(s)$ of energy alone. Let $A_{R_0}$ be the pole coefficient of
this pair, with the normalization preceding the common factor $1/(8sH)$.
Its $z$-cut finite part is
\begin{equation}
 F_{R_0+W,z}=E_{R_0+W}(s)
 -\frac{2s(-s^5-5s^4+18s^3+122s^2+151s-45)}{3(s^2+2s-3)}
 +A_{R_0}\ln\frac{s-1}{s}.
\end{equation}
The rational subtraction is the complete lower-basepoint constant.
The $W$ part has zero lower constant and zero soft pole. Finally, the
previously derived rational-sector matching shift converts this expression
to the common $(xy)^\eta$ convention. This order of operations preserves
all finite endpoint terms before combining the other real sectors.

\subsection{Explicit combined rational and \texorpdfstring{$W$}{W} primitive}
Put $r=\sqrt{s}>1$. All GPLs below have basepoint $r$ and endpoint $u$. The following finite sum specifies the complete combined primitive, including its rational and elementary terms.
\begin{equation}
 \mathcal F_{R_0+W}(u)=\sum_{j=0}^{37}c_j(u)\,\mathcal G_j(u)+\sum_{j=1}^3 d_j(u)\mathcal L_j(u)
 +K\left[G(0)\{G(0,s)-G(0,1)\}-G(0)^3/3\right].
\end{equation}
Here $K=32s(2s+1)/(s-1)^2$. The two weight-three GPL terms already occur in the coefficient table, so the last term contains only their shuffle products. Write $q_\pm=(r^2-3\pm\sqrt{r^4-10r^2+9})/2$ and $a_\pm=1\pm\sqrt{1-r^2}$, $b_\pm=-1\pm\sqrt{1-r^2}$.
Define the polynomial abbreviations before listing the coefficients.
\begin{align}
P_{1}={}&540- 1380 s- 537 s^{5}- 48 s^{6}+12 s^{7}+610 s^{4}\\
&+2334 s^{3}+12701 s^{2}.
\end{align}
\begin{align}
P_{2}={}&252- 1128 s- 609 s^{6}- 24 s^{7}+12 s^{8}+784 s^{5}\\
&+1872 s^{4}+2705 s^{2}+14424 s^{3}.
\end{align}
\begin{align}
P_{3}={}&5- s^{4}- s^{5}- 49 s+20 s^{2}+26 s^{3}+107 u- 109 s u\\
&- 47 s^{2} u+25 s^{3} u.
\end{align}
\begin{align}
P_{4}={}&90- 42 s- 42 s^{2}- 6 s^{3}+117 u+s^{4} u+24 s^{3} u\\
&+130 s^{2} u+208 s u.
\end{align}
\begin{align}
P_{5}={}&2268- 227127 s^{3}- 184189 s^{4}- 36221 s^{5}- 13500 s\\
&- 2693 s^{8}- 2367 s^{2}- 597 s^{9}- 319 s^{7}+12 s^{11}\\
&+36 s^{10}+6201 s^{6}.
\end{align}
\begin{align}
P_{6}={}&-648- 5122 s^{4}- 3759 s^{5}- 1589 s^{3}- 1246 s^{6}\\
&- 187 s^{7}+2 s^{8}+3 s^{9}+972 s+2358 s^{2}.
\end{align}
\begin{align}
P_{7}={}&1755- 10593 s^{2}- 6172 s- 4076 s^{3}- 55 s^{4}+5 s^{6}\\
&+128 s^{5}.
\end{align}
\begin{align}
P_{8}={}&-162+s^{9}- 2476 s^{3}- 1980 s^{2}- 1409 s^{4}- 689 s^{5}\\
&- 286 s^{6}- 70 s^{7}- 3 s^{8}+162 s.
\end{align}
\begin{align}
P_{9}={}&-162+s^{9}- 2692 s^{3}- 2153 s^{4}- 1008 s^{2}- 989 s^{5}\\
&- 322 s^{6}- 70 s^{7}- 3 s^{8}+486 s.
\end{align}
\begin{align}
\mathcal G_{0}&=1,\\
\mathcal G_{1}&=G(0;u),\\
\mathcal G_{2}&=G(1;u),\\
\mathcal G_{3}&=G(r^{2};u).
\end{align}
\begin{align}
c_{0}={}&\frac{P_{1} u}{18 s \left(s - 1\right)^{2}}\\
&- \frac{u^{2} \left(12 s^{5} + 261 s^{4} - 700 s^{3} + 1027 s^{2} + 2040 s - 396\right)}{36 s \left(s - 1\right)^{2}}\\
&+\frac{u^{3} \left(12 s^{4} + 24 s^{3} - 79 s^{2} + 787 s + 48\right)}{54 \left(s - 1\right)^{2}}\\
&- \frac{u^{4} \left(24 s^{2} + 185 s - 4\right)}{24 \left(s - 1\right)^{2}}\\
&+\frac{2 s u^{5} \left(s + 7\right)}{15 \left(s - 1\right)^{2}}\\
&- \frac{P_{1}}{18 u \left(s - 1\right)^{2}}\\
&- \frac{2 s^{6} \left(s + 7\right)}{15 u^{5} \left(s - 1\right)^{2}}\\
&- \frac{s^{3} \left(12 s^{4} + 24 s^{3} - 79 s^{2} + 787 s + 48\right)}{54 u^{3} \left(s - 1\right)^{2}}\\
&+\frac{s^{4} \left(24 s^{2} + 185 s - 4\right)}{24 u^{4} \left(s - 1\right)^{2}}\\
&+\frac{s \left(12 s^{5} + 261 s^{4} - 700 s^{3} + 1027 s^{2} + 2040 s - 396\right)}{36 u^{2} \left(s - 1\right)^{2}}\\
&+\frac{8 s \left(s + 5\right) \left(s + u\right)}{\left(s + 3\right) \left(s + u^{2} + 2 u\right)}\\
&+\frac{4 s \left(- s + u\right) \left(s + 5\right) \left(s^{3} - 3 s^{2} - 9 s + 1\right)}{3 \left(s - 1\right) \left(s + u^{2} - 2 u\right)}.
\end{align}
\begin{align}
c_{1}={}&\frac{P_{2}}{18 s \left(s - 1\right)^{2}}\\
&+\frac{4 u \left(s^{5} - 13 s^{4} + 12 s^{3} - 123 s^{2} + 45 s - 9\right)}{3 s \left(s - 1\right)^{2}}\\
&+\frac{u^{2} \left(3 s^{4} + 4 s^{3} + 169 s^{2} + 12 s - 6\right)}{3 s \left(s - 1\right)^{2}}\\
&- \frac{4 s u^{3} \left(s + 26\right)}{9 \left(s - 1\right)^{2}}\\
&+\frac{5 s u^{4}}{6 \left(s - 1\right)^{2}}\\
&+\frac{4 \left(s^{5} - 13 s^{4} + 12 s^{3} - 123 s^{2} + 45 s - 9\right)}{3 u \left(s - 1\right)^{2}}\\
&+\frac{5 s^{5}}{6 u^{4} \left(s - 1\right)^{2}}\\
&- \frac{4 P_{3} s}{3 \left(s - 1\right) \left(s + u^{2} - 2 u\right)}\\
&- \frac{4 s^{4} \left(s + 26\right)}{9 u^{3} \left(s - 1\right)^{2}}\\
&+\frac{s \left(3 s^{4} + 4 s^{3} + 169 s^{2} + 12 s - 6\right)}{3 u^{2} \left(s - 1\right)^{2}}\\
&+\frac{4 P_{4} s}{3 \left(s + 3\right)^{2} \left(s + u^{2} + 2 u\right)}\\
&+\frac{4 s \left(- s + 2 u\right) \left(s + 5\right) \left(s^{3} - 3 s^{2} - 9 s + 1\right)}{3 \left(s + u^{2} - 2 u\right)^{2}}\\
&+\frac{8 s \left(s - 1\right) \left(s + 5\right) \left(s + 2 u\right)}{\left(s + 3\right) \left(s + u^{2} + 2 u\right)^{2}}.
\end{align}
\begin{align}
c_{2}={}&- \frac{P_{5}}{18 s \left(s - 1\right)^{3} \left(s + 3\right)^{2}}\\
&+\frac{s u^{4}}{2 \left(s - 1\right)^{2}}\\
&+\frac{4 u \left(7 s^{6} - 31 s^{5} - 82 s^{4} - 138 s^{3} - 415 s^{2} - 30 s - 3\right)}{3 s \left(s - 1\right)^{3}}\\
&- \frac{u^{2} \left(17 s^{5} - 39 s^{4} - 37 s^{3} - 237 s^{2} - 102 s + 18\right)}{3 s \left(s - 1\right)^{3}}\\
&+\frac{4 s u^{3} \left(3 s - 10\right)}{9 \left(s - 1\right)^{2}}\\
&+\frac{s^{5}}{2 u^{4} \left(s - 1\right)^{2}}\\
&+\frac{4 \left(7 s^{6} - 31 s^{5} - 82 s^{4} - 138 s^{3} - 415 s^{2} - 30 s - 3\right)}{3 u \left(s - 1\right)^{3}}\\
&- \frac{4 P_{3} s}{3 \left(s - 1\right) \left(s + u^{2} - 2 u\right)}\\
&- \frac{s \left(17 s^{5} - 39 s^{4} - 37 s^{3} - 237 s^{2} - 102 s + 18\right)}{3 u^{2} \left(s - 1\right)^{3}}\\
&+\frac{4 s^{4} \left(3 s - 10\right)}{9 u^{3} \left(s - 1\right)^{2}}\\
&+\frac{4 s \left(- s + 2 u\right) \left(s + 5\right) \left(s^{3} - 3 s^{2} - 9 s + 1\right)}{3 \left(s + u^{2} - 2 u\right)^{2}}.
\end{align}
\begin{align}
c_{3}={}&\frac{P_{5}}{18 s \left(s - 1\right)^{3} \left(s + 3\right)^{2}}\\
&- \frac{s u^{4}}{2 \left(s - 1\right)^{2}}\\
&- \frac{4 u \left(7 s^{6} - 31 s^{5} - 82 s^{4} - 138 s^{3} - 415 s^{2} - 30 s - 3\right)}{3 s \left(s - 1\right)^{3}}\\
&+\frac{u^{2} \left(17 s^{5} - 39 s^{4} - 37 s^{3} - 237 s^{2} - 102 s + 18\right)}{3 s \left(s - 1\right)^{3}}\\
&- \frac{4 s u^{3} \left(3 s - 10\right)}{9 \left(s - 1\right)^{2}}\\
&- \frac{4 \left(7 s^{6} - 31 s^{5} - 82 s^{4} - 138 s^{3} - 415 s^{2} - 30 s - 3\right)}{3 u \left(s - 1\right)^{3}}\\
&- \frac{s^{5}}{2 u^{4} \left(s - 1\right)^{2}}\\
&- \frac{4 s^{4} \left(3 s - 10\right)}{9 u^{3} \left(s - 1\right)^{2}}\\
&+\frac{s \left(17 s^{5} - 39 s^{4} - 37 s^{3} - 237 s^{2} - 102 s + 18\right)}{3 u^{2} \left(s - 1\right)^{3}}\\
&+\frac{4 P_{3} s}{3 \left(s - 1\right) \left(s + u^{2} - 2 u\right)}\\
&- \frac{4 s \left(- s + 2 u\right) \left(s + 5\right) \left(s^{3} - 3 s^{2} - 9 s + 1\right)}{3 \left(s + u^{2} - 2 u\right)^{2}}.
\end{align}
\begin{align}
\mathcal G_{4}&=G(b_{-},0;u),\\
\mathcal G_{5}&=G(b_{+},0;u),\\
\mathcal G_{6}&=G(a_{-},0;u),\\
\mathcal G_{7}&=G(a_{+},0;u).
\end{align}
\begin{equation}
c_{4}=\frac{4 s \left(s^{5} + 9 s^{4} + 116 s^{3} + 516 s^{2} + 307 s - 693\right)}{3 \left(s - 1\right) \left(s + 3\right)^{3}}.
\end{equation}
\begin{equation}
c_{5}=\frac{4 s \left(s^{5} + 9 s^{4} + 116 s^{3} + 516 s^{2} + 307 s - 693\right)}{3 \left(s - 1\right) \left(s + 3\right)^{3}}.
\end{equation}
\begin{equation}
c_{6}=- \frac{4 \left(s^{6} - s^{5} - 33 s^{4} + 75 s^{3} + 18 s^{2} + 24 s - 12\right)}{3 s \left(s - 1\right)^{2}}.
\end{equation}
\begin{equation}
c_{7}=- \frac{4 \left(s^{6} - s^{5} - 33 s^{4} + 75 s^{3} + 18 s^{2} + 24 s - 12\right)}{3 s \left(s - 1\right)^{2}}.
\end{equation}
\begin{align}
\mathcal G_{8}&=G(q_{-},0;u),\\
\mathcal G_{9}&=G(q_{+},0;u),\\
\mathcal G_{10}&=G(a_{-},1;u),\\
\mathcal G_{11}&=G(a_{+},1;u).
\end{align}
\begin{equation}
c_{8}=- \frac{128 s \left(s + 5\right)}{3 \left(s - 1\right)}.
\end{equation}
\begin{equation}
c_{9}=- \frac{128 s \left(s + 5\right)}{3 \left(s - 1\right)}.
\end{equation}
\begin{equation}
c_{10}=- \frac{4 \left(s^{6} - s^{5} - 33 s^{4} + 75 s^{3} + 18 s^{2} + 24 s - 12\right)}{3 s \left(s - 1\right)^{2}}.
\end{equation}
\begin{equation}
c_{11}=- \frac{4 \left(s^{6} - s^{5} - 33 s^{4} + 75 s^{3} + 18 s^{2} + 24 s - 12\right)}{3 s \left(s - 1\right)^{2}}.
\end{equation}
\begin{align}
\mathcal G_{12}&=G(0,1;u),\\
\mathcal G_{13}&=G(q_{-},1;u),\\
\mathcal G_{14}&=G(q_{+},1;u),\\
\mathcal G_{15}&=G(a_{-},r^{2};u).
\end{align}
\begin{equation}
c_{12}=- \frac{4 \left(s^{8} - 14 s^{7} - 93 s^{6} + 30 s^{5} + 403 s^{4} - 2002 s^{3} - 8747 s^{2} - 8334 s + 324\right)}{3 \left(s - 1\right)^{2} \left(s + 3\right)^{3}}.
\end{equation}
\begin{equation}
c_{13}=\frac{128 s \left(s + 5\right)}{3 \left(s - 1\right)}.
\end{equation}
\begin{equation}
c_{14}=\frac{128 s \left(s + 5\right)}{3 \left(s - 1\right)}.
\end{equation}
\begin{equation}
c_{15}=\frac{4 \left(s^{6} - s^{5} - 33 s^{4} + 75 s^{3} + 18 s^{2} + 24 s - 12\right)}{3 s \left(s - 1\right)^{2}}.
\end{equation}
\begin{align}
\mathcal G_{16}&=G(a_{+},r^{2};u),\\
\mathcal G_{17}&=G(0,r^{2};u),\\
\mathcal G_{18}&=G(q_{-},r^{2};u),\\
\mathcal G_{19}&=G(q_{+},r^{2};u).
\end{align}
\begin{equation}
c_{16}=\frac{4 \left(s^{6} - s^{5} - 33 s^{4} + 75 s^{3} + 18 s^{2} + 24 s - 12\right)}{3 s \left(s - 1\right)^{2}}.
\end{equation}
\begin{align}
c_{17}={}&- \frac{4 P_{6}}{3 s \left(s - 1\right)^{2} \left(s + 3\right)^{3}}.
\end{align}
\begin{equation}
c_{18}=- \frac{128 s \left(s + 5\right)}{3 \left(s - 1\right)}.
\end{equation}
\begin{equation}
c_{19}=- \frac{128 s \left(s + 5\right)}{3 \left(s - 1\right)}.
\end{equation}
\begin{align}
\mathcal G_{20}&=G(- r^{2};u),\\
\mathcal G_{21}&=G(b_{-},- r^{2};u),\\
\mathcal G_{22}&=G(b_{+},- r^{2};u),\\
\mathcal G_{23}&=G(0,- r^{2};u).
\end{align}
\begin{align}
c_{20}={}&- \frac{2 P_{7} s}{3 \left(s - 1\right)^{3} \left(s + 3\right)^{2}}\\
&- \frac{4 s u \left(s^{3} + 45 s^{2} - 405 s + 43\right)}{3 \left(s - 1\right)^{3}}\\
&+\frac{2 s u^{2} \left(s^{2} - 26 s + 63\right)}{3 \left(s - 1\right)^{3}}\\
&- \frac{4 P_{4} s}{3 \left(s + 3\right)^{2} \left(s + u^{2} + 2 u\right)}\\
&- \frac{4 s^{2} \left(s^{3} + 45 s^{2} - 405 s + 43\right)}{3 u \left(s - 1\right)^{3}}\\
&+\frac{2 s^{3} \left(s^{2} - 26 s + 63\right)}{3 u^{2} \left(s - 1\right)^{3}}\\
&- \frac{8 s \left(s - 1\right) \left(s + 5\right) \left(s + 2 u\right)}{\left(s + 3\right) \left(s + u^{2} + 2 u\right)^{2}}.
\end{align}
\begin{equation}
c_{21}=- \frac{4 s \left(s^{5} + 9 s^{4} + 116 s^{3} + 516 s^{2} + 307 s - 693\right)}{3 \left(s - 1\right) \left(s + 3\right)^{3}}.
\end{equation}
\begin{equation}
c_{22}=- \frac{4 s \left(s^{5} + 9 s^{4} + 116 s^{3} + 516 s^{2} + 307 s - 693\right)}{3 \left(s - 1\right) \left(s + 3\right)^{3}}.
\end{equation}
\begin{equation}
c_{23}=\frac{4 s \left(s^{4} - 2 s^{3} - 678 s^{2} + 150 s - 431\right)}{3 \left(s - 1\right)^{3}}.
\end{equation}
\begin{align}
\mathcal G_{24}&=G(1,- r^{2};u),\\
\mathcal G_{25}&=G(r^{2},- r^{2};u),\\
\mathcal G_{26}&=G(q_{-},- r^{2};u),\\
\mathcal G_{27}&=G(q_{+},- r^{2};u).
\end{align}
\begin{equation}
c_{24}=- \frac{256 s \left(s^{6} - 53 s^{4} - 231 s^{3} - 388 s^{2} - 253 s - 36\right)}{3 \left(s - 1\right)^{3} \left(s + 3\right)^{3}}.
\end{equation}
\begin{equation}
c_{25}=- \frac{256 s \left(s^{6} - 53 s^{4} - 231 s^{3} - 388 s^{2} - 253 s - 36\right)}{3 \left(s - 1\right)^{3} \left(s + 3\right)^{3}}.
\end{equation}
\begin{equation}
c_{26}=\frac{256 s \left(s + 5\right)}{3 \left(s - 1\right)}.
\end{equation}
\begin{equation}
c_{27}=\frac{256 s \left(s + 5\right)}{3 \left(s - 1\right)}.
\end{equation}
\begin{align}
\mathcal G_{28}&=G(-1;u),\\
\mathcal G_{29}&=G(b_{-},-1;u),\\
\mathcal G_{30}&=G(b_{+},-1;u),\\
\mathcal G_{31}&=G(0,-1;u).
\end{align}
\begin{align}
c_{28}={}&\frac{2 P_{7} s}{3 \left(s - 1\right)^{3} \left(s + 3\right)^{2}}\\
&+\frac{4 s u \left(s^{3} + 45 s^{2} - 405 s + 43\right)}{3 \left(s - 1\right)^{3}}\\
&- \frac{2 s u^{2} \left(s^{2} - 26 s + 63\right)}{3 \left(s - 1\right)^{3}}\\
&- \frac{2 s^{3} \left(s^{2} - 26 s + 63\right)}{3 u^{2} \left(s - 1\right)^{3}}\\
&+\frac{4 P_{4} s}{3 \left(s + 3\right)^{2} \left(s + u^{2} + 2 u\right)}\\
&+\frac{4 s^{2} \left(s^{3} + 45 s^{2} - 405 s + 43\right)}{3 u \left(s - 1\right)^{3}}\\
&+\frac{8 s \left(s - 1\right) \left(s + 5\right) \left(s + 2 u\right)}{\left(s + 3\right) \left(s + u^{2} + 2 u\right)^{2}}.
\end{align}
\begin{equation}
c_{29}=\frac{4 s \left(s^{5} + 9 s^{4} + 116 s^{3} + 516 s^{2} + 307 s - 693\right)}{3 \left(s - 1\right) \left(s + 3\right)^{3}}.
\end{equation}
\begin{equation}
c_{30}=\frac{4 s \left(s^{5} + 9 s^{4} + 116 s^{3} + 516 s^{2} + 307 s - 693\right)}{3 \left(s - 1\right) \left(s + 3\right)^{3}}.
\end{equation}
\begin{equation}
c_{31}=- \frac{4 s \left(s^{4} - 2 s^{3} - 678 s^{2} + 150 s - 431\right)}{3 \left(s - 1\right)^{3}}.
\end{equation}
\begin{align}
\mathcal G_{32}&=G(q_{-},-1;u),\\
\mathcal G_{33}&=G(q_{+},-1;u),\\
\mathcal G_{34}&=G(1,-1;u),\\
\mathcal G_{35}&=G(r^{2},-1;u).
\end{align}
\begin{equation}
c_{32}=- \frac{256 s \left(s + 5\right)}{3 \left(s - 1\right)}.
\end{equation}
\begin{equation}
c_{33}=- \frac{256 s \left(s + 5\right)}{3 \left(s - 1\right)}.
\end{equation}
\begin{equation}
c_{34}=\frac{256 s \left(s^{6} - 53 s^{4} - 231 s^{3} - 388 s^{2} - 253 s - 36\right)}{3 \left(s - 1\right)^{3} \left(s + 3\right)^{3}}.
\end{equation}
\begin{equation}
c_{35}=\frac{256 s \left(s^{6} - 53 s^{4} - 231 s^{3} - 388 s^{2} - 253 s - 36\right)}{3 \left(s - 1\right)^{3} \left(s + 3\right)^{3}}.
\end{equation}
\begin{align}
\mathcal G_{36}&=G(0,0,1;u),\\
\mathcal G_{37}&=G(0,0,r^{2};u).
\end{align}
\begin{equation}
c_{36}=\frac{32 s \left(2 s + 1\right)}{\left(s - 1\right)^{2}}.
\end{equation}
\begin{equation}
c_{37}=- \frac{32 s \left(2 s + 1\right)}{\left(s - 1\right)^{2}}.
\end{equation}
\begin{equation}
\mathcal L_{1}=G(0;u)G(0;u).
\end{equation}
\begin{align}
d_{1}={}&- \frac{2 \left(s^{5} - 22 s^{4} + 22 s^{3} - 398 s^{2} + 73 s + 36\right)}{3 \left(s - 1\right)^{2}}\\
&- \frac{16 s u \left(s + 1\right)}{\left(s - 1\right)^{2}}\\
&- \frac{8 s^{3}}{u^{2} \left(s - 1\right)^{2}}\\
&+\frac{8 s u^{2}}{\left(s - 1\right)^{2}}\\
&- \frac{32 s \left(s u - 3 s + 3 u - 1\right)}{\left(- s + u^{2}\right) \left(s - 1\right)}\\
&+\frac{16 s^{2} \left(s + 1\right)}{u \left(s - 1\right)^{2}}.
\end{align}
\begin{equation}
\mathcal L_{2}=G(0;u)G(1;u).
\end{equation}
\begin{align}
d_{2}={}&\frac{8 P_{8}}{3 s \left(s - 1\right)^{2} \left(s + 3\right)^{3}}\\
&- \frac{16 s u \left(s + 1\right)}{\left(s - 1\right)^{2}}\\
&- \frac{8 s^{3}}{u^{2} \left(s - 1\right)^{2}}\\
&+\frac{8 s u^{2}}{\left(s - 1\right)^{2}}\\
&- \frac{32 s \left(s u - 3 s + 3 u - 1\right)}{\left(- s + u^{2}\right) \left(s - 1\right)}\\
&+\frac{16 s^{2} \left(s + 1\right)}{u \left(s - 1\right)^{2}}.
\end{align}
\begin{equation}
\mathcal L_{3}=G(0;u)G(r^{2};u).
\end{equation}
\begin{align}
d_{3}={}&\frac{8 P_{9}}{3 s \left(s - 1\right)^{2} \left(s + 3\right)^{3}}\\
&+\frac{16 s u \left(s + 1\right)}{\left(s - 1\right)^{2}}\\
&- \frac{8 s u^{2}}{\left(s - 1\right)^{2}}\\
&+\frac{8 s^{3}}{u^{2} \left(s - 1\right)^{2}}\\
&- \frac{16 s^{2} \left(s + 1\right)}{u \left(s - 1\right)^{2}}\\
&+\frac{32 s \left(s u - 3 s + 3 u - 1\right)}{\left(- s + u^{2}\right) \left(s - 1\right)}.
\end{align}
All coefficients follow by collecting the original two primitives and applying the displayed shuffle identities. The generating script checks their exact coefficient reconstruction before this table is used. Endpoint constants and the common regulator conversion are applied to the combined primitive as before.

\paragraph{The grouped upper endpoint.} The energy-only weight-two result can be grouped into seven coefficient
blocks. Put $\bar c_j=c_j(s)$ and $\bar d_j=d_j(s)$, using the explicit
coefficient table, and define
\begin{align}
 D_b&=\sum_{\epsilon=\pm}
       (d_{b_\epsilon,0}-d_{b_\epsilon,-s}+d_{b_\epsilon,-1}),\\
 D_a&=\sum_{\epsilon=\pm}
       (d_{a_\epsilon,0}+d_{a_\epsilon,1}-d_{a_\epsilon,s}),\\
 D_q&=\sum_{\epsilon=\pm}
       (d_{q_\epsilon,0}-d_{q_\epsilon,1}+d_{q_\epsilon,s}
                      -2d_{q_\epsilon,-s}+2d_{q_\epsilon,-1}).
\end{align}
Then all 30 weight-two terms have the compact form
\begin{align}
 E^{(2)}_{R_0+W}={}&\bar c_4D_b+\bar c_6D_a+\bar c_8D_q
       +\bar c_{12}d_{0,1}+\bar c_{17}d_{0,s}\\
 &+\bar c_{23}(d_{0,-s}-d_{0,-1})
       +\bar c_{24}(d_{1,-s}+d_{s,-s}-d_{1,-1}-d_{s,-1}).
\end{align}
The coefficient identities are rational identities for arbitrary $s>1$.
The elementary part is
\begin{align}
 E^{(\mathrm{el})}_{R_0+W}={}&\bar c_0+\bar c_1\ell
 +\bar c_2\ln[r(s-1)]+\bar c_{20}\ln\frac{2r}{s+1}\\
 &+\bar d_1\ell^2+\bar d_2\ell\ln(r+1)
                         -\bar d_3\ell\ln[r(r-1)].
\end{align}
Here the cancellations $\bar c_3=-\bar c_2$ and
$\bar c_{28}=-\bar c_{20}$ have been performed before displaying the
answer. The complete constant upper endpoint is
$E^{(\mathrm{el})}_{R_0+W}+E^{(2)}_{R_0+W}
+\mathcal F^{(3)}_{R_0+W}(s)$.


\subsection{Cancellation between the four real sectors}
Section~\ref{sec:comp-real-kernel-collection} has already removed the
$b_\pm$ pole family directly from the common integrand. The following
identities also recover that cancellation from the separately integrated
sectors and reduce the remaining functions.
The common primitive becomes shorter when the rational, $W$,
energy-logarithm and Bose-restoring contributions are combined before
reducing individual functions. Denote these primitives by
$F_{R_0}$, $F_W$, $F_{\log}$ and $F_D$, respectively. The factor
$\ln s$ multiplying the energy-logarithm primitive is included in
$F_{\log}$. All $G$ functions in this subsection have the same
basepoint $r=\sqrt{s}$ and the same final argument $u$.
The subscript $R_0$ denotes the rational kernel. Taking its integrated
finite part with the same cutoff $z>\zeta$ as in the sector assembly gives
\begin{equation}
 F_{R_0}^z=F_{\mathrm{rat}}^z.
\end{equation}
This is a change of notation for the same finite part, with no
additional matching shift or normalization factor.

Consider the pair of letters
\begin{equation}
 b_\pm=-1\pm\sqrt{1-s},\qquad
 \sqrt{1-s}=\ii\sqrt{s-1},\qquad s>1.
\end{equation}
Their complete contribution to the rational primitive is
\begin{align}
 \left.F_{R_0}\right|_{b_\pm}
 &=c_b(s)\sum_{b=b_+,b_-}
       \bigl[G(b,0;u)-G(b,-s;u)+G(b,-1;u)\bigr],\\
 c_b(s)&=\frac{4s\left(s^5+9s^4+116s^3+516s^2+307s-693\right)}
                   {3(s-1)(s+3)^3}.
\end{align}
The independently integrated Bose-restoring primitive has exactly the
opposite six coefficients:
\begin{equation}
 \left.F_D\right|_{b_\pm}=-\left.F_{R_0}\right|_{b_\pm},
 \qquad \left.F_W\right|_{b_\pm}=\left.F_{\log}\right|_{b_\pm}=0.
 \label{eq:comp-real-four-b-cancel}
\end{equation}
Thus neither letter is needed in the combined primitive. This is a
cancellation of identical functions with identical paths, so it introduces
no change of branches or boundary constants.

The letters $r$ and $-r$ cancel between the rational and $W$ sectors.
Define
\begin{equation}
 Q_c(u)=G(c,0;u)+G(c,1;u)-G(c,s;u).
\end{equation}
Their coefficients are
\begin{align}
 \left.F_{R_0}\right|_{\pm r}
 &=-\frac{16(r-1)^2}{r+1}Q_r(u)
       +\frac{16(r+1)^2}{r-1}Q_{-r}(u),\\
 \left.F_W\right|_{\pm r}&=-\left.F_{R_0}\right|_{\pm r},
 \qquad \left.F_{\log}\right|_{\pm r}=\left.F_D\right|_{\pm r}=0.
\end{align}
Combining all four sectors therefore removes twelve weight-two words
with these four first letters. The remaining alphabet is
\begin{equation}
 \left\{0,1,-1,s,-s,1+\sqrt{1-s},1-\sqrt{1-s},q_+,q_-\right\},
 \qquad q_\pm=\frac{s-3\pm\sqrt{(s-1)(s-9)}}2.
\end{equation}
The $q_\pm$ pair has the alternative representation derived above.
It is not necessary to expand every remaining $G$ function to make
these cancellations explicit.

There are 44, 19, 14 and 28 nonzero GPL terms in the four separate
primitive dictionaries. After collecting their coefficients, 37 distinct
words remain. Three shuffle reductions leave 34 words and three products
of weight-one functions, together with the rational term. The coefficient
collection uses $H=s-1$ and $r^2=s$ before identifying equal letters.
Both the change of variables in each coefficient and the shuffle
reconstruction are checked as exact algebraic identities.
This collection precedes the finite endpoint operation. The lower-limit
rational constants, the moving-boundary finite shifts, and the soft
contribution must still be combined to obtain a compact total cross
section.

\paragraph{The finite endpoint after collection.}
Let $F(u)=F_{R_0}(u)+F_W(u)+F_{\log}(u)+F_D(u)$ and
$b=s-u$. The collected primitive has only two rational coefficients
with a pole at $b=0$: its constant coefficient $c_{\varnothing}$ and
its coefficient $c_0$ multiplying $G(0;u)=\ln(u/\sqrt{s})$.
Their residues are
\begin{align}
 c_{\varnothing}(s-b)&=\frac{d_0(s)}b+O(1),
 &d_0(s)&=\frac{8s(s^3+6s^2-s-2)}3,\\
 c_0(s-b)&=\frac{d_1(s)}b+O(1),
 &d_1(s)&=-\frac{32s^2(s-5)}{3(s-1)}.
\end{align}
All other function coefficients are regular at $u=s$.
The product with the first-order expansion of $G(0;u)$ is essential:
\begin{equation}
 G(0;s-b)=\tfrac12\ln s-\frac bs+O(b^2).
\end{equation}
It contributes the finite term
\begin{equation}
 T_{\mathrm{end}}(s)=-\frac{d_1(s)}s
                 =\frac{32s(s-5)}{3(s-1)}.
 \label{eq:comp-real-four-endpoint-Taylor}
\end{equation}
Thus substituting $u=s$ in the regular coefficients alone would miss a
finite contribution. The coefficient of the double pole in the
integrand is the coefficient of $1/b$ in its primitive:
\begin{equation}
 D_4(s)=d_0(s)+\tfrac12d_1(s)\ln s.
 \label{eq:comp-real-four-double-pole}
\end{equation}

At the lower limit, every GPL vanishes. The coefficients of the
products of two weight-one functions have at most a simple pole at
$u=\sqrt{s}$, whereas the products vanish quadratically, so these terms
also give zero. The rational term has the finite limit
\begin{equation}
 F(\sqrt{s})=-\frac{2(s+1)(s^4-3s^3-41s^2+27s-8)}{3(s-1)}.
 \label{eq:comp-real-four-lower-limit}
\end{equation}
This combines the four lower limits before evaluating any polylogarithm.

To state the remaining change of cutoff, define $A_4(s)$ by the
integrand expansion
\begin{equation}
 F'(s-b)=\frac{D_4(s)}{b^2}+\frac{A_4(s)}b+O(\ln b).
\end{equation}
In particular the primitive contains $D_4/b-A_4\ln b$.
Writing $\mathcal E_4(s)$ for its finite coefficient, including
\eqref{eq:comp-real-four-endpoint-Taylor}, the four-sector finite part
in the original $z$ convention is
\begin{equation}
 F_4^z=\mathcal E_4(s)-F(\sqrt{s})
       +A_4(s)\ln\frac{s-1}{s}-\frac{D_4(s)}{s(s-1)}.
 \label{eq:comp-real-four-zfinite}
\end{equation}
Indeed, $z=(s-1)b/s-b^2/s^2+O(b^3)$ gives
\begin{equation}
 \frac1b=\frac{s-1}{sz}-\frac1{s(s-1)}+O(z),
 \qquad \ln b=\ln z+\ln\frac{s}{s-1}+O(z).
\end{equation}
The second-order boundary relation supplies the last finite term in
\eqref{eq:comp-real-four-zfinite}. The finite shifts from the $z$
convention to the common real-emission convention are subsequently
added as in \eqref{eq:comp-matching-short}.
\paragraph{Removing a conjugate pair at the endpoint.}
The remaining letters $a_\pm=1\pm\ii\sqrt{s-1}$ occur only in the
combination
\begin{equation}
 \mathcal A_{\mathrm{pair}}(u)=\sum_{a=a_+,a_-}
       \bigl[G(a,0;u)+G(a,1;u)-G(a,s;u)\bigr].
\end{equation}
Their common coefficient in the four-sector primitive is
\begin{equation}
 \gamma_a(s)=-\frac{32s(s^2-3s-1)}{(s-1)^2}\ln s
 -\frac{4(s^5+s^4-25s^3+47s^2-12s+12)}{3s(s-1)}.
\end{equation}
It is shorter to integrate the summed logarithmic derivative than to
evaluate the six complex-letter GPLs separately. Put
\begin{equation}
 B(u)=\frac{u(u-1)}{s-u},\qquad
 P(u)=u^2-2u+s=(u-a_+)(u-a_-).
\end{equation}
For $\sqrt{s}<u<s$, $B(u)>0$ and $B(\sqrt{s})=1$.
The second-letter combination is exactly $\ln B(u)$, so
\begin{equation}
 \mathcal A_{\mathrm{pair}}'(u)=\frac{P'(u)}{P(u)}\ln B(u).
\end{equation}
The factorization $P(u)=(s-u)(1+B(u))$ gives
\begin{equation}
 \frac{P'(u)}{P(u)}=\frac{B'(u)}{1+B(u)}+\frac1{u-s}.
\end{equation}
Since the derivative of
$\ln B\ln(1+B)+\operatorname{Li}_2(-B)$ is
$\ln B\,B'/(1+B)$, the primitive is
\begin{equation}
 \begin{aligned}
 \mathcal A_{\mathrm{pair}}(u)={}&\ln B(u)\ln(1+B(u))+\operatorname{Li}_2(-B(u))
                      +\frac{\pi^2}{12}\\
 &+G(s,0;u)+G(s,1;u)-G(s,s;u).
 \end{aligned}
\end{equation}
The constant follows from $\mathcal A_{\mathrm{pair}}(\sqrt{s})=0$ and
$\operatorname{Li}_2(-1)=-\pi^2/12$. All arguments in the explicit
dilogarithm term are negative real, fixing its branch throughout the
integration interval.

The separate terms contain logarithms at the upper endpoint, but their
sum has a finite limit. With $r=\sqrt{s}$ and $b=s-u$, the required
expansions are
\begin{align}
 \ln B\ln(1+B)+\operatorname{Li}_2(-B)
 &=\tfrac12\ln^2\!\left(\frac{s(s-1)}b\right)
      -\frac{\pi^2}{6}+O(b|\ln b|),\\
 G(s,0;s-b)&=\ln r\ln\frac b{s-r}
           +\operatorname{Li}_2(1-1/r)+O(b),\\
 G(s,1;s-b)&=\ln(r+1)\ln\frac b{s-r}
           +\operatorname{Li}_2\!\left(\frac r{r+1}\right)+O(b),\\
 G(s,s;s-b)&=\tfrac12\ln^2\frac b{s-r}.
\end{align}
Both powers of $\ln b$ cancel exactly. The limit first takes the form
\begin{equation}
 \mathcal A_{\mathrm{pair}}(s)=\operatorname{Li}_2(1-1/r)
 +\operatorname{Li}_2\!\left(\frac r{r+1}\right)
 +\tfrac12\ln^2[r(r+1)]-\frac{\pi^2}{12}.
\end{equation}
The reflection and inversion identities of the dilogarithm then give
\begin{equation}
 \boxed{\mathcal A_{\mathrm{pair}}(s)=\frac{\pi^2}{4}
       +\operatorname{Li}_2(-1/\sqrt{s})
       -\operatorname{Li}_2(1/\sqrt{s})
       +\frac12\ln s\ln(s-1).}
 \label{eq:comp-real-a-pair-endpoint}
\end{equation}
For example, use
$\operatorname{Li}_2(r/(r+1))=-\operatorname{Li}_2(-r)
-\tfrac12\ln^2(r+1)$ and
$\operatorname{Li}_2(-r)+\operatorname{Li}_2(-1/r)
=-\pi^2/6-\tfrac12\ln^2r$, together with the reflection identity
for $\operatorname{Li}_2(1-1/r)$.
The final two dilogarithm arguments lie in $(-1,1)$ for every $s>1$.
Thus the complete contribution is
$\gamma_a(s)\mathcal A_{\mathrm{pair}}(s)$, and the six complex-letter endpoint values
are unnecessary.

\subsection{The matched real contribution}
The four real sectors and their finite matching shifts can now be combined into a formula with two dilogarithms and one trilogarithm. Define
\begin{equation}
 \mathcal R_4(s)=F_{\mathrm{rat}}^z+F_W^z+F_{\log}^z+F_D^z
                    +\Delta_{\mathrm{rat}}+\Delta_{\log}+\Delta_D.
\end{equation}
Here $s>1$, $H=s-1$ and $r=\sqrt{s}$. Use the real logarithms
\begin{equation}
 L=\ln s,\qquad M=\ln(s-1),\qquad J=\ln\frac{s+1}{2},
\end{equation}
and the following three special-function combinations:
\begin{align}
 D(s)&=\operatorname{Li}_2(1/s),\\
 V(s)&=\operatorname{Li}_2\!\left(\frac2{s+1}\right)+\ln2\,J,\\
 T(s)&=\zeta(3)-\operatorname{Li}_3(1/s).
\end{align}
Both dilogarithm arguments and the trilogarithm argument lie strictly between zero and one. The complete matched expression is
\begin{equation}
\begin{aligned}
 \mathcal R_4(s)={}&A_0+A_1L+A_2L^2+A_3L^3
                  +\pi^2(B_0+B_1L)\\
 &+(C_0+C_1L)D+KT+NV+EJ^2+FJ\\
 &+M(H_0+H_1L+H_2L^2+H_3J).
\end{aligned}
\label{eq:comp-real-four-matched-compact}
\end{equation}
Every coefficient is a rational function of $s$. Explicitly,
\begin{equation}
A_0=\frac{\begin{aligned}&30 s^{7} - 198 s^{6} - 785 s^{5} + 4079 s^{4}\\&- 22927 s^{3} + 22309 s^{2} - 10692 s + 4428\end{aligned}}{108 s H}.
\end{equation}
\begin{equation}
A_1=\frac{\begin{aligned}&39 s^{7} - 333 s^{6} + 174 s^{5} + 10076 s^{4}\\&+6798 s^{3} - 116 s^{2} + 222 s + 36\end{aligned}}{9 s H^{3}}.
\end{equation}
\begin{equation}
A_2=\frac{\begin{aligned}&- 5 s^{7} + 77 s^{6} + 102 s^{5} - 630 s^{4}\\&- 641 s^{3} - 151 s^{2} + 132 s - 36\end{aligned}}{3 s H^{3}}.
\end{equation}
\begin{equation}
A_3=- \frac{8 s \left(3 s^{2} - 25 s - 2\right)}{3 H^{2}}.
\end{equation}
\begin{equation}
B_0=\frac{2 \left(s^{4} - 17 s^{3} + 25 s^{2} + 219 s - 12\right)}{9 H}.
\end{equation}
\begin{equation}
B_1=- \frac{8 s \left(5 s + 1\right)}{3 H^{2}}.
\end{equation}
\begin{equation}
C_0=- \frac{4 \left(s^{4} - 17 s^{3} - 39 s^{2} - 101 s - 12\right)}{3 H}.
\end{equation}
\begin{equation}
C_1=\frac{16 s}{s - 1}.
\end{equation}
\begin{equation}
K=\frac{32 s \left(2 s + 1\right)}{H^{2}}.
\end{equation}
\begin{equation}
N=- \frac{256 s \left(s + 5\right)}{3 H}.
\end{equation}
\begin{equation}
E=- \frac{128 s \left(2 s^{3} - 5 s^{2} - 24 s + 3\right)}{3 H^{3}}.
\end{equation}
\begin{equation}
F=\frac{256 s \left(s + 1\right) \left(s^{2} - 19 s - 4\right)}{3 H^{3}}.
\end{equation}
\begin{equation}
H_0=\frac{\begin{aligned}&- 2 s^{7} + 2 s^{6} + 64 s^{5} - 260 s^{4}\\&- 314 s^{3} - 2638 s^{2} + 220 s - 144\end{aligned}}{3 s H}.
\end{equation}
\begin{equation}
H_1=\frac{\begin{aligned}&- 8 s^{7} + 80 s^{6} + 24 s^{5} - 1800 s^{4}\\&- 2336 s^{3} - 2248 s^{2} + 240 s - 96\end{aligned}}{3 s H^{3}}.
\end{equation}
\begin{equation}
H_2=- \frac{16 s \left(2 s^{2} - 9 s - 2\right)}{H^{2}}.
\end{equation}
\begin{equation}
H_3=- \frac{256 s \left(s^{3} - 19 s^{2} - 21 s - 9\right)}{3 H^{3}}.
\end{equation}
This is the unnormalized real finite contribution. Its relation to the remainder in \eqref{eq:comp-simplified-master} is
\begin{equation}
 \mathcal R_{\mathrm{rem}}(X)=\mathcal R_4(1+X)+2F_{UW}^{\eta}(X).
\end{equation}
Thus \eqref{eq:comp-real-four-matched-compact} replaces the four separate endpoint evaluations and their three matching shifts. The UV term is already included with the virtual contribution in that master formula. The remaining $UW$ and finite soft terms are separate.
\paragraph{Reduction of the function basis.} After the endpoint identities above, real reflection and inversion relations put the dilogarithm arguments in $(0,1)$. Two five-term identities reduce their number further:
\begin{align}
 \operatorname{Li}_2\!\left(\frac2{r+1}\right)+\operatorname{Li}_2\!\left(\frac{r+1}{2r}\right)
 &=\operatorname{Li}_2(1/r)+\operatorname{Li}_2\!\left(\frac1{r+1}\right)+\frac{\pi^2}{12}-\frac12\ln^22+\ln2\ln\frac{r+1}{r}.
\end{align}
The same identity with $r$ replaced by $s$ reduces the half-argument function in the matching shift. For the second identity, define $\mathcal L(z)=\operatorname{Li}_2(z)+\tfrac12\ln z\ln(1-z)$ for $0<z<1$. Then
\begin{equation}
 \mathcal L\!\left(\frac{r+1}{r^2+1}\right)-\mathcal L\!\left(\frac{r-1}{r(r+1)}\right)
 =\mathcal L(1/r)+\mathcal L\!\left(\frac1{r^2+1}\right)-\mathcal L\!\left(\frac1{r+1}\right).
\end{equation}
Both identities follow by differentiation using $\dd\operatorname{Li}_2(z)/\dd z=-\ln(1-z)/z$, with the integration constant fixed as $r\to\infty$. The remaining square-root arguments occur in a common-coefficient difference. Landen transformation and duplication give
\begin{equation}
 \operatorname{Li}_2\!\left(\frac1{r+1}\right)-\operatorname{Li}_2(1/r)
 =-\frac12\operatorname{Li}_2(1/r^2)-\frac12\ln^2\!\left(\frac{r+1}{r}\right).
\end{equation}
After this substitution, all coefficients of $\ln(r+1)$ cancel exactly when $\ln(r-1)=\ln(s-1)-\ln(r+1)$ is used. The result therefore has no square-root argument. Adding the matching shift also cancels the entire coefficient of $\operatorname{Li}_2(1/(s+1))$.
The matching shifts are reduced with the explicit identities
\begin{align}
 \operatorname{Li}_2(1-s)&=D-LM+\tfrac12L^2-\frac{\pi^2}{6},\\
 \operatorname{Li}_2\!\left(-\frac{(s-1)^2}{4s}\right)
 &=4\operatorname{Li}_2\!\left(\frac2{s+1}\right)-2D-2\operatorname{Li}_2\!\left(\frac1{s+1}\right)\\
 &\quad+J^2+2J(L+\ln2)-\tfrac32L^2+2M(L-2J)-\frac{\pi^2}{6}-\ln^22.
\end{align}
For the second line, put $v=(s-1)/(s+1)$, apply Landen transformation to the negative argument, and then use $\operatorname{Li}_2(v^2)=2\operatorname{Li}_2(v)+2\operatorname{Li}_2(-v)$. Reflection and the preceding half-argument identity give the displayed expression. All intermediate arguments used in these identities are real and their logarithms have the signs fixed by $s>1$.
Equivalently, denote the two right-hand sides by $Q_1(s)$ and $Q_2(s)$. Differentiation and the positive-factor logarithm identities for $s>1$ give
\begin{equation}
 Q_1'(s)=-\frac{\ln s}{s-1},\qquad
 Q_2'(s)=-\frac{2(s+1)}{s(s-1)}\ln\frac{s+1}{2\sqrt s}.
\end{equation}
These are precisely the derivatives of the two left-hand sides. At $s\to1^+$ both sides vanish. Here the coefficient of $M=\ln(s-1)$ vanishes at least linearly, and $\operatorname{Li}_2(1/2)=\pi^2/12-\ln^22/2$ fixes the remaining constant. Equality of derivatives and this boundary value establish the identities on the connected domain $s>1$.

\subsection{The integrated UW contribution}
Put $H=X>0$ and use $Y$ and $q$ defined at the beginning of this chapter. The completed $UW$ integral can be expressed in the same functions as the virtual contribution. Define the real logarithms
\begin{equation} L_0=\ln X,\qquad L_1=\ln(1+X),\qquad L_2=\ln(1+X/2).\end{equation}
In the following formula every $G$ has endpoint $X$. Its rational letters have kernels $\dd z/(z-a)$, and the letter $r$ has kernel $\dd z/\sqrt{z(z+4)}$, with positive square root for $z>0$. Define
\begin{equation} \mathcal U(X)=G(-4,r,-1;X)+G(r,-1,-1;X)-G(r,0,-1;X).\end{equation}
The unnormalized finite contribution is
\begin{equation}\begin{aligned}
 F_{UW}^{\eta}(X)={}&a_0+a_1L_1+a_LL_0+a_{11}L_1^2+a_{10}L_0L_1\\
 &+a_D\operatorname{Li}_2(-X)+a_2L_2+a_{D2}\operatorname{Li}_2(-X/2)\\
 &+a_{20}(L_0L_2+L_2^2/4)+(a_R+2KL_0)G(r,-1;X)\\
 &+K\mathcal U(X)+a_{011}G(0,-1,-1;X)+a_{RR}G(r,r,-1;X).
\end{aligned}\label{eq:comp-uw-integrated-compact}\end{equation}
Its contribution to $C_1$ is $F_{UW}^{\eta}/[4X(1+X)]$, including the second Bose-related sector. The coefficients are as follows.
\begin{equation}
a_0=\frac{\begin{aligned}&497 X^{6} - 1647 X^{5} + 46 X^{4} + 51270 X^{3}\\&+18924 X^{2} - 46080 X - 18432\end{aligned}}{216 X^{3}}.
\end{equation}
\begin{equation}
a_1=\frac{\begin{aligned}&- 24 X^{8} + 24 X^{7} + 1548 X^{6} - 7712 X^{5}\\&- 91001 X^{4} - 208525 X^{3} - 177316 X^{2} - 44544 X\\&+6144\end{aligned}}{18 X^{4} \left(X + 4\right)}.
\end{equation}
\begin{equation}
a_L=\frac{\begin{aligned}&- 12 X^{6} - 24 X^{5} + 48 X^{4} + 2912 X^{3}\\&+6656 X^{2} + 5376 X + 1536\end{aligned}}{9 X^{3}}.
\end{equation}
\begin{equation}
a_{11}=\frac{\begin{aligned}&X^{8} - 3 X^{7} - 106 X^{6} + 210 X^{5}\\&+3762 X^{4} + 8930 X^{3} + 7656 X^{2} + 1664 X\\&-512\end{aligned}}{3 X^{4} \left(X + 4\right)}.
\end{equation}
\begin{equation}
a_{10}=\frac{\begin{aligned}&4 X^{8} - 12 X^{7} - 352 X^{6} + 144 X^{5}\\&+8160 X^{4} + 21248 X^{3} + 19200 X^{2} + 3584 X\\&-2048\end{aligned}}{3 X^{4} \left(X + 4\right)}.
\end{equation}
\begin{equation}
a_D=\frac{\begin{aligned}&2 X^{6} - 6 X^{5} - 224 X^{4} - 632 X^{3}\\&+352 X^{2} + 1920 X + 1152\end{aligned}}{3 X^{2} \left(X + 4\right)}.
\end{equation}
\begin{equation}
a_2=- \frac{128 \left(X + 1\right) \left(X + 2\right) \left(X^{2} - 17 X - 22\right)}{3 X^{3}}.
\end{equation}
\begin{equation}
a_{D2}=\frac{128 \left(X + 1\right) \left(X + 6\right)}{3 X}.
\end{equation}
\begin{equation}
a_{20}=\frac{256 \left(X + 1\right) \left(X^{3} - 5 X^{2} - 28 X - 24\right)}{3 X^{3}}.
\end{equation}
\begin{equation}
a_R=\frac{\begin{aligned}&- 24 X^{5} Y - 72 X^{4} Y + 528 X^{3} Y + 2112 X^{2} Y\\&+768 X Y - 768 Y\end{aligned}}{X^{3} \left(X + 4\right)}.
\end{equation}
\begin{equation}
K=\frac{\begin{aligned}&24 X^{5} Y + 72 X^{4} Y - 528 X^{3} Y - 2496 X^{2} Y\\&- 2688 X Y - 768 Y\end{aligned}}{X^{3} \left(X + 4\right)}.
\end{equation}
\begin{equation}
a_{011}=- \frac{16 \left(X + 1\right) \left(7 X + 3\right)}{X^{2}}.
\end{equation}
\begin{equation}
a_{RR}=- \frac{24 \left(X + 1\right) \left(X^{2} - 2 X - 12\right)}{X^{2}}.
\end{equation}
\paragraph{Integration and endpoint constants.} Starting from the eight-function kernel above, rational integration by parts gives
\begin{equation} \Phi_{UW}(t)=\frac{\dd}{\dd t}\sum_i Q_i(t)f_i(t)+\sum_i T_i(t)f_i(t).\end{equation}
Here $Q_i$ are rational functions, and $T_i$ have only simple poles. The identity follows by substituting the derivatives of the eight functions and matching their rational coefficients. At each endpoint the negative powers in the Laurent expansion of the sum $\sum_iQ_if_i$ cancel. Its constant term must include products of the poles of $Q_i$ with the Taylor coefficients of $f_i$. Thus the integrated boundary is the difference of the two constant terms, not the difference of the constant coefficients of the rational functions alone.
The simple-pole terms integrate to GPLs with endpoint $t=1$.
Set $X=(1-q)^2/q$ and $\delta=1-q$, with $0<q<1$. The moving letters are
\begin{equation}
 a=-\frac{q}{(1-q)^2},\qquad c=-\frac{2q}{(1-q)^2},\qquad
 r_-=-\frac{q}{1-q},\qquad r_+=\frac1{1-q}.
\end{equation}
They lie outside the interval $[0,1]$. Together with the fixed letters
$0,1$, they specify the endpoint functions whose constants are needed.

For a word $w$, define $g_w(q)$ as the constant coefficient in the
$t\to1^-$ expansion of $G(w;t)$ in powers of $\ln(1-t)$.
The endpoint is regularized at fixed $q$ before taking $q\to1$.
The remaining terms vanish as $(1-t)$ times powers of its logarithm,
so this constant-term operation respects shuffle products.
In particular, $g_{\varnothing}=1$, $g_0=g_1=0$, and all nonempty
pure-zero or pure-one words have value zero. The nonzero fixed-letter
values through weight three are
\begin{align}
 g_{0,1}&=-\zeta(2),&g_{1,0}&=\zeta(2),\\
 g_{0,0,1}&=-\zeta(3),&g_{0,1,0}&=2\zeta(3),\\
 g_{1,0,0}&=-\zeta(3),&g_{0,1,1}&=\zeta(3),\\
 g_{1,0,1}&=-2\zeta(3),&g_{1,1,0}&=\zeta(3).
\end{align}
Three convergent integrals fix this table:
\begin{align}
 G(0,1;1)&=\int_0^1\frac{\ln(1-t)}t\,\dd t=-\zeta(2),\\
 G(0,0,1;1)&=-\operatorname{Li}_3(1)=-\zeta(3),\\
 G(0,1,1;1)&=\frac12\int_0^1\frac{\ln^2(1-t)}t\,\dd t
 =\frac12\int_0^1\frac{\ln^2u}{1-u}\,\dd u=\zeta(3).
\end{align}
The first two follow from the convergent polylogarithm series.
For the third, expanding $1/(1-u)=\sum_{n=0}^{\infty}u^n$ and using
$\int_0^1u^n\ln^2u\,\dd u=2/(n+1)^3$ gives the stated value.
The remaining entries follow from the regularized shuffle relations
\begin{align}
 0&=g_{0,1}+g_{1,0},\\
 0&=2g_{0,0,1}+g_{0,1,0},\\
 0&=g_{1,0,0}+g_{0,1,0}+g_{0,0,1},\\
 0&=g_{1,0,1}+2g_{0,1,1},\\
 0&=g_{0,1,1}+g_{1,0,1}+g_{1,1,0}.
\end{align}
These are the products $g_0g_1$, $g_0g_{0,1}$, $g_1g_{0,0}$,
$g_1g_{0,1}$ and $g_0g_{1,1}$, respectively.

The words containing a moving letter have zero boundary value as
$q\to1$. To establish this before interchanging limits, first use
shuffle products to remove initial letters $1$ and terminal letters $0$.
The factors $g_1$ and $g_0$ vanish. Every surviving word therefore
starts with a letter other than $1$ and ends with a letter other than
$0$, and still contains each moving letter of the original word.
Its ordered integral is absolutely convergent at both endpoints.
For $0<\delta<1/4$, each moving letter $\alpha$ satisfies
\begin{equation}
 \left|\frac1{t-\alpha(q)}\right|\leq\frac2{|\alpha(q)|},
 \qquad 0\leq t\leq1.
\end{equation}
Replacing the moving kernels by these bounds leaves only the fixed
kernels $1/t$, $1/(1-t)$ and the constant kernel in the ordered integral.
The endpoint conditions just established make this integral finite.
Consequently, for each surviving word $w$,
\begin{equation}
 |G(w;1)|\leq C_w
 \prod_{j\,\mathrm{moving}}\frac2{|\alpha_j(q)|},
 \qquad C_w<\infty,
\end{equation}
where $C_w$ is independent of $q$. The factors $1/a$ and $1/c$
are $O(\delta^2)$, while $1/r_\pm=O(\delta)$. The bound therefore
tends to zero, and so does the finite shuffle combination defining
$g_w$. This proves the moving-letter boundary value through every
weight required here. It applies to each special function before
multiplication by possibly singular rational energy coefficients.

Differentiating the endpoint GPLs by the total-differential recursion
of Section~\ref{sec:comp-virtual-energy-basis} now gives logarithmic
differentials in a fixed alphabet with the explicitly fixed constants
above. Add the complete Hermite boundary difference, including its
Laurent--Taylor cross terms, before returning to the $X$ kernels.
The result contains eighteen words with rational coefficients in $X$
and $Y$. Exact substitution of the kernel maps reproduces every
coefficient of the fixed-alphabet expression.
\paragraph{Collection of the functions.} The shuffle identity gives
\begin{equation}G(0,r,-1;X)=L_0G(r,-1;X)-G(r,0,-1;X)-G(r,-1,0;X).\end{equation}
The five root words occur with coefficients proportional to $1,2,1,2,1$, in the order $(-4,r,-1)$, $(0,r,-1)$, $(r,-1,-1)$, $(r,-1,0)$, $(r,0,-1)$. Applying this identity reduces their sum to $K[\mathcal U+2L_0G(r,-1;X)]$. For the rational letters, use
\begin{align}G(-1,0;X)&=L_0L_1+\operatorname{Li}_2(-X),\\
G(-2,0;X)&=L_0L_2+\operatorname{Li}_2(-X/2),\\
G(-j,-j;X)&=\tfrac12\ln^2(1+X/j),\qquad j=1,2.\end{align}
These are shuffle identities and the defining dilogarithm integral, so no continuation across a cut is needed. All kernels and logarithms in the final expression are real for $X>0$.

\section{The collected total in one function basis}
\label{sec:comp-total-energy-collected}
All contributions can now be combined in the same energy variable. This section gives the complete coefficient of $\alpha_{\mathrm T}^3/m^2$, including the virtual correction and the integrated two-photon final state. It is obtained from the independently evaluated sectors above. Define $X=s_{\mathrm{phys}}/m^2-1>0$ and $Y=\sqrt{X/(X+4)}$. Use the GPL kernels and real branches in Section~\ref{sec:comp-virtual-energy-basis}, including the removable combinations containing the letter $1$. The complete result is
\begin{equation}
\sigma_{\mathrm{inc}}^{\mathrm{NLO}}=\frac{\pi\alpha_{\mathrm T}^2}{m^2}B(X)+\frac{\alpha_{\mathrm T}^3}{m^2}C_1(X)+\frac1{m^2}O(\alpha_{\mathrm T}^4),
\end{equation}
where $B$ is given in \eqref{eq:comp-simplified-born} and
\begin{equation}\boxed{C_1(X)=\sum_{j=0}^{22}c_j(X,Y)F_j(X).}\label{eq:comp-total-energy-collected}\end{equation}
The functions and their full coefficients follow. No unevaluated phase-space or Feynman-parameter integral occurs in this expression.
\begin{equation}
 F_{0}=1,\qquad c_{0}=\frac{- 273 X^{3} + 982 X^{2} + 2960 X + 1744}{24 X^{2} \left(X + 1\right)^{2}}.
\end{equation}
\begin{equation}
 F_{1}=G(-1;X),\qquad c_{1}=\frac{37 X^{4} - 54 X^{3} - 339 X^{2} - 428 X - 184}{4 X^{3} \left(X + 1\right)^{2}}.
\end{equation}
\begin{equation}
 F_{2}=G(0;X),\qquad c_{2}=\frac{- 14 X^{4} - 17 X^{3} + 17 X^{2} + 22 X + 8}{2 X \left(X - 1\right) \left(X + 1\right)^{3}}.
\end{equation}
\begin{equation}
 F_{3}=G(-1,-1;X),\qquad c_{3}=\frac{\begin{aligned}&- 4 X^{6} - 35 X^{5} + 31 X^{4} + 755 X^{3}\\&+1765 X^{2} + 1506 X + 440\end{aligned}}{X^{3} \left(X + 1\right)^{2} \left(X + 4\right)}.
\end{equation}
\begin{equation}
 F_{4}=G(-1,0;X),\qquad c_{4}=\frac{X^{4} + 7 X^{3} + X^{2} - 3 X - 2}{X^{3} \left(X + 1\right)^{2}}.
\end{equation}
\begin{equation}
 F_{5}=G(-2,-1;X),\qquad c_{5}=\frac{\begin{aligned}&- 4 X^{5} - 104 X^{4} - 584 X^{3} - 1264 X^{2}\\&- 1152 X - 384\end{aligned}}{X^{3} \left(X + 1\right)^{2} \left(X + 4\right)}.
\end{equation}
\begin{equation}
 F_{6}=G(0,-1;X),\qquad c_{6}=\frac{\begin{aligned}&X^{6} + 8 X^{5} - 17 X^{4} - 210 X^{3}\\&- 448 X^{2} - 352 X - 96\end{aligned}}{X^{3} \left(X + 1\right)^{2} \left(X + 4\right)}.
\end{equation}
\begin{equation}
 F_{7}=G(1,0;X),\qquad c_{7}=\frac{- X^{2} + X - 2}{X^{3}}.
\end{equation}
\begin{equation}
 F_{8}=G(r,-1;X),\qquad c_{8}=\frac{Y \left(3 X^{4} + 18 X^{3} + 44 X^{2} - 8 X - 64\right)}{X^{4}}.
\end{equation}
\begin{equation}
 F_{9}=G(-1,-1,-1;X),\qquad c_{9}=\frac{- 2 X^{2} + 32 X + 46}{X^{3}}.
\end{equation}
\begin{equation}
 F_{10}=G(-1,-1,0;X),\qquad c_{10}=\frac{\left(2 X + 3\right) \left(4 X - 8\right)}{X^{3}}.
\end{equation}
\begin{equation}
 F_{11}=G(-1,-2,-1;X),\qquad c_{11}=\frac{8 X^{2} - 32 X - 48}{X^{3}}.
\end{equation}
\begin{equation}
 F_{12}=G(-1,0,-1;X),\qquad c_{12}=\frac{4 X^{2} - 14 X - 24}{X^{3}}.
\end{equation}
\begin{equation}
 F_{13}=G(-4,r,-1;X),\qquad c_{13}=- \frac{6 Y \left(X^{3} - 12 X - 8\right)}{X^{4}}.
\end{equation}
\begin{equation}
 F_{14}=G(0,-1,-1;X),\qquad c_{14}=\frac{- 5 X^{2} - 32 X + 8}{X^{3}}.
\end{equation}
\begin{equation}
 F_{15}=G(0,-1,0;X),\qquad c_{15}=\frac{2 X^{2} + 4 X - 12}{X^{3}}.
\end{equation}
\begin{equation}
 F_{16}=G(0,1,0;X),\qquad c_{16}=\frac{2 X^{2} + 4 X - 12}{X^{3}}.
\end{equation}
\begin{equation}
 F_{17}=G(0,r,-1;X),\qquad c_{17}=- \frac{3 Y \left(X - 2\right) \left(X + 4\right)}{X^{3}}.
\end{equation}
\begin{equation}
 F_{18}=G(r,-1,-1;X),\qquad c_{18}=\frac{3 Y \left(5 X^{4} + 14 X^{3} - 96 X^{2} - 352 X - 128\right)}{X^{4} \left(X + 4\right)}.
\end{equation}
\begin{equation}
 F_{19}=G(r,-1,0;X),\qquad c_{19}=- \frac{6 Y \left(X - 2\right) \left(X + 4\right)}{X^{3}}.
\end{equation}
\begin{equation}
 F_{20}=G(r,-2,-1;X),\qquad c_{20}=- \frac{16 Y \left(X^{4} + 2 X^{3} - 24 X^{2} - 80 X - 48\right)}{X^{4} \left(X + 4\right)}.
\end{equation}
\begin{equation}
 F_{21}=G(r,0,-1;X),\qquad c_{21}=- \frac{8 Y \left(X^{4} + 3 X^{3} - 18 X^{2} - 68 X - 24\right)}{X^{4} \left(X + 4\right)}.
\end{equation}
\begin{equation}
 F_{22}=G(r,r,-1;X),\qquad c_{22}=\frac{9 X^{2} - 24}{X^{3}}.
\end{equation}
\paragraph{Collection of the real terms.} Put $L=\ln(1+X)$, $M=\ln X$, and $J=\ln(1+X/2)$. The classical functions in the four-sector real expression satisfy
\begin{align}
\operatorname{Li}_2\!\left(\frac1{1+X}\right)&=\zeta(2)+G(-1,0;X)-G(-1,-1;X),\\
\operatorname{Li}_2\!\left(\frac2{2+X}\right)+\ln2\,J&=\zeta(2)+G(-2,0;X)-G(-2,-2;X),\\
\zeta(3)-\operatorname{Li}_3\!\left(\frac1{1+X}\right)&=\zeta(2)G(-1;X)+G(-1,-1,0;X)-G(-1,-1,-1;X).
\end{align}
Each identity follows by differentiating with respect to $X$. The constants are fixed as $X\to0^+$, where the dilogarithms and trilogarithm take their values at one. Every nonconstant GPL on these right-hand sides vanishes in this limit. In particular, the logarithmic lower endpoint in $G(-1,0;X)$ and $G(-2,0;X)$ is integrable because $\int_0^X|\ln t|\,\mathrm{d}t=O(X[1+|\ln X|])$. The remaining products of $L$, $M$, and $J$ are expanded with the shuffle identity, counting each interleaving even when repeated letters give the same word. All classical arguments above lie in $(0,1)$ for $X>0$, so this conversion introduces no continuation across a cut.
\paragraph{Assembly.} Substitute these identities, the completed $UW$ expression, and \eqref{eq:comp-ck-energy-compact} into \eqref{eq:comp-simplified-master}. The finite virtual coefficient is its constant term in $\ell_\lambda=\ln(\lambda_{\mathrm{phys}}/m)$ after the previously proved infrared cancellation. Taking this coefficient does not set the photon mass to $m$. The finite conversion term remains $-A(1+X)[\ln X+\tfrac12\ln(1+X)]$. Expanding this product gives the last coefficients needed in the sum above.
There are five contributions to this collection: $\widehat C_{V+\mathrm{UV}}$, $F_{UW}^{\eta}/[4X(1+X)]$, $\mathcal R_4(1+X)/[8X(1+X)]$, $-\mathcal C_K(1+X)/(1+X)$, and the finite conversion term. The factor $1/4$ for $UW$ follows from $\mathcal R_{\mathrm{rem}}=\mathcal R_4+2F_{UW}^{\eta}$, so its Bose-related second copy is already included. Expand each contribution in the same words and add their rational coefficients. This is an algebraic reconstruction and does not require a claim that the functions are linearly independent. The infrared coefficient used in the last contribution is \eqref{eq:comp-ir-short}, with $L^2=2G(-1,-1;X)$ and $\mathcal G(X)=G(r,-1;X)$.
The coefficients of five words vanish identically in the complete sum. In particular the coefficient of $G(-1,r,-1;X)$ from the finite soft term cancels its coefficient from the infrared conversion term. All coefficients proportional to $\pi^2$ outside the GPLs also cancel. The remaining coefficients are the rational functions listed above.
\paragraph{The apparent point $X=1$.} The factor $1/(X-1)$ multiplies $G(0;X)=\ln X$, so it has a removable limit. The product involving this coefficient tends to $+1$. The two GPLs containing the rational letter $+1$ are also real and finite:
\begin{align}
G(1,0;X)&=\ln X\ln|1-X|+\operatorname{Re}\operatorname{Li}_2(X),\\
G(0,1,0;X)&=2\operatorname{Re}\operatorname{Li}_3(X)-\ln X\operatorname{Re}\operatorname{Li}_2(X).
\end{align}
Their derivatives are $\ln X/(X-1)$ and $G(1,0;X)/X$. At $X=1$ their values are $\zeta(2)$ and $2\zeta(3)$. Thus the apparent inner pole does not require a contour choice for the physical answer.

\subsection{A product form of the total}
\label{sec:comp-total-energy-product}
\begingroup
\def\comptonproductlabel{eq:comp-total-energy-product}
Products of elementary logarithms can make the collected result shorter without introducing additional polylogarithms. Define $L=\ln(1+X)$ and $M=\ln X$. The two GPLs with the same coefficient combine into
\begin{equation}
\mathcal P(X)=G(0,-1,0;X)+G(0,1,0;X)=\tfrac12\left[\operatorname{Re}\operatorname{Li}_3(X^2)-M\operatorname{Re}\operatorname{Li}_2(X^2)\right].
\end{equation}
For the derivation, first take $0<X<1$. The shuffle relation $G(0)G(0,a)=2G(0,0,a)+G(0,a,0)$, with $G(0,a;X)=-\operatorname{Li}_2(X/a)$ and $G(0,0,a;X)=-\operatorname{Li}_3(X/a)$, gives
\begin{equation}
 G(0,a,0;X)=2\operatorname{Li}_3(X/a)-M\operatorname{Li}_2(X/a),\qquad a=\pm1.
\end{equation}
Summing the two signs and using $\operatorname{Li}_n(z)+\operatorname{Li}_n(-z)=2^{1-n}\operatorname{Li}_n(z^2)$ gives the stated result. Both shuffle multiplicities and this duplication factor are needed for its factor $1/2$. For $X>1$, continue all functions on the same edge of the principal cut. On either edge, the imaginary parts of $\operatorname{Li}_2(X)$ and $\operatorname{Li}_3(X)$ are respectively $\pm\pi\ln X$ and $\pm\pi\ln^2X/2$, with the same sign. These follow by integrating their derivatives from $X=1$, and they cancel in $2\operatorname{Li}_3(X)-M\operatorname{Li}_2(X)$. The analogous cancellation holds for the paired expression at argument $X^2$. Thus both edges give the same real function and the continuation adds no branch constant. The following equivalent form uses only real polylogarithms with arguments in $(0,1]$:
\begin{equation}
\mathcal P(X)=\begin{cases}\tfrac12[\operatorname{Li}_3(X^2)-M\operatorname{Li}_2(X^2)],&0<X\leq1,\\ \tfrac12[\operatorname{Li}_3(X^{-2})+M\operatorname{Li}_2(X^{-2})]+\tfrac13M^3+\tfrac16\pi^2M,&X\geq1.\end{cases}
\end{equation}
To derive the second line, put $t=\ln X>0$, $D(t)=\operatorname{Li}_2(e^{-2t})$ and $T(t)=\operatorname{Li}_3(e^{-2t})$. Their derivatives are $D^{\prime}(t)=2\ln(1-e^{-2t})$ and $T^{\prime}(t)=-2D(t)$. The identity $\ln(e^{2t}-1)=2t+\ln(1-e^{-2t})$ then gives
\begin{align}
\operatorname{Re}\operatorname{Li}_2(e^{2t})&=\frac{\pi^2}{3}-2t^2-D(t),\\
\operatorname{Re}\operatorname{Li}_3(e^{2t})&=T(t)-\frac43t^3+\frac{2\pi^2}{3}t.
\end{align}
Indeed, the derivatives of the two sides are respectively $-2\ln(e^{2t}-1)$ and $2\operatorname{Re}\operatorname{Li}_2(e^{2t})$. Write $\zeta(n)=\sum_{k=1}^{\infty}k^{-n}$ for the Riemann zeta function at integers $n>1$. Their values at $t=0$ agree because $\operatorname{Li}_2(1)=\pi^2/6$ and $\operatorname{Li}_3(1)=\zeta(3)$. These derivative and boundary conditions fix both identities and give $\mathcal P(1)=\zeta(3)/2$.
The sole apparent rational pole at $X=1$ can also be removed before evaluation. Define the continuous elementary function
\begin{equation}
R_{\log}(X)=\begin{cases}\ln X/(X-1),&X\ne1,\\1,&X=1,\end{cases}\qquad R_{\log}(1+h)=1-\frac h2+\frac{h^2}{3}+O(h^3).
\end{equation}
The series follows by dividing the Taylor series of $\ln(1+h)$ by $h$, which converges for $|h|<1$. Extracting the common power $X^{-3}$ gives the exact total coefficient as
\begin{equation}
C_1(X)=\frac{1}{X^3}\left[R_{\log}(X)+\mathcal E(X,L,M)+\sum_{j=1}^{13}b_j(X,Y,L,M)\mathcal F_j(X)\right].
\label{\comptonproductlabel}\end{equation}
The elementary part is
\begin{equation}
\mathcal E=e_{0}L^{3}+e_{1}L^{2}M+e_{2}L^{2}+e_{3}LM+e_{4}L+e_{5}M+e_{6}.
\end{equation}
\begin{equation}
e_{0}=- \frac{X^{2} - 16 X - 23}{3}.
\end{equation}
\begin{equation}
e_{1}=4 X^{2} - 2 X - 12.
\end{equation}
\begin{equation}
e_{2}=- 2 X^{3} - \frac{11 X^{2}}{2} + \frac{133 X}{2} + 36 + \frac{96}{X + 4} - \frac{6}{X + 1} + \frac{1}{\left(X + 1\right)^{2}}.
\end{equation}
\begin{equation}
e_{3}=X^{2} + 5 X - 10 + \frac{12}{X + 1} - \frac{4}{\left(X + 1\right)^{2}}.
\end{equation}
\begin{equation}
e_{4}=\frac{37 X^{2}}{4} - 32 X - 30 - \frac{15}{X + 1} - \frac{1}{\left(X + 1\right)^{2}}.
\end{equation}
\begin{equation}
e_{5}=- \frac{14 X^{5} + 31 X^{4} + 14 X^{3} - 6 X^{2} - 8 X - 2}{2 \left(X + 1\right)^{3}}.
\end{equation}
\begin{equation}
e_{6}=- \frac{X \left(273 X^{3} - 982 X^{2} - 2960 X - 1744\right)}{24 \left(X + 1\right)^{2}}.
\end{equation}
The thirteen remaining functions and their coefficients are
\begin{equation}
\mathcal F_{1}=G(0,-1;X).
\end{equation}
\begin{equation}
b_{1}=X^{3} + X^{2} - 43 X + 6 - \frac{96}{X + 4} - \frac{6}{X + 1} + \frac{2}{\left(X + 1\right)^{2}}- 2 L X \left(2 X + 5\right).
\end{equation}
\begin{equation}
\mathcal F_{2}=G(-2,-1;X).
\end{equation}
\begin{equation}
b_{2}=- 4 X^{2} - 80 X - 68 - \frac{128}{X + 4} + \frac{8}{X + 1} - \frac{4}{\left(X + 1\right)^{2}}.
\end{equation}
\begin{equation}
\mathcal F_{3}=G(1,0;X),\qquad b_{3}=- X^{2} + X - 2.
\end{equation}
\begin{equation}
\mathcal F_{4}=G(r,-1;X).
\end{equation}
\begin{equation}
b_{4}=Y \left(3 X^{3} + 18 X^{2} + 44 X - 8 - \frac{64}{X}\right)- 3 M Y \left(X - 2\right) \left(X + 4\right).
\end{equation}
\begin{equation}
\mathcal F_{5}=G(0,-1,-1;X),\qquad b_{5}=- 5 X^{2} - 8 X + 32.
\end{equation}
\begin{equation}
\mathcal F_{6}=G(-1,-2,-1;X),\qquad b_{6}=8 \left(X^{2} - 4 X - 6\right).
\end{equation}
\begin{equation}
\mathcal F_{7}=G(-4,r,-1;X),\qquad b_{7}=Y \left(- 6 X^{2} + 72 + \frac{48}{X}\right).
\end{equation}
\begin{equation}
\mathcal F_{8}=\mathcal P(X),\qquad b_{8}=2 X^{2} + 4 X - 12.
\end{equation}
\begin{equation}
\mathcal F_{9}=G(r,0,-1;X),\qquad b_{9}=Y \left(- 5 X^{2} + 14 X + 88 + \frac{48}{X + 4} + \frac{48}{X}\right).
\end{equation}
\begin{equation}
\mathcal F_{10}=G(r,-1,0;X),\qquad b_{10}=- 3 Y \left(X - 2\right) \left(X + 4\right).
\end{equation}
\begin{equation}
\mathcal F_{11}=G(r,-1,-1;X),\qquad b_{11}=Y \left(15 X^{2} - 18 X - 216 - \frac{96}{X + 4} - \frac{96}{X}\right).
\end{equation}
\begin{equation}
\mathcal F_{12}=G(r,-2,-1;X),\qquad b_{12}=Y \left(- 16 X^{2} + 32 X + 256 + \frac{64}{X + 4} + \frac{192}{X}\right).
\end{equation}
\begin{equation}
\mathcal F_{13}=G(r,r,-1;X),\qquad b_{13}=9 X^{2} - 24.
\end{equation}

\endgroup
This form follows from the complete coefficient table using only the identities
\begin{align}
G(-1,0)&=LM-G(0,-1),\\
G(-1,0,-1)&=L\,G(0,-1)-2G(0,-1,-1),\\
G(-1,-1,0)&=\tfrac12ML^2-L\,G(0,-1)+G(0,-1,-1),\\
G(0,r,-1)&=M\,G(r,-1)-G(r,0,-1)-G(r,-1,0).
\end{align}
The common endpoint is $X$. In addition, repeated letters $-1$ give powers of $L$ divided by factorials. Expanding every product back by the shuffle identity reproduces all twenty-three original coefficients exactly. The two words combined into $\mathcal P$ have the same coefficient $c_{15}=c_{16}$, so this last grouping replaces their sum without another factor of two.
The extraction of $R_{\log}$ is fixed by the coefficient $c_2$ of $M$ in the collected table. Direct polynomial division gives
\begin{equation}
 X^3c_2(X,Y)=\frac{1}{X-1}+e_5(X),\qquad X^3c_2(X,Y)M=R_{\log}(X)+e_5(X)M.
\end{equation}
The second equality at $X=1$ means its continuous limit. All other coefficients are multiplied by $X^3$ before the common factor $X^{-3}$ is placed outside the bracket. Thus this extraction changes neither the normalization of $C_1$ nor the coefficient of its removable term. Every displayed rational coefficient is regular at $X=1$, as follows by exact polynomial cancellation. The function $G(1,0;X)=\zeta(2)-\operatorname{Li}_2(1-X)$ is also explicitly real and regular for every $X>0$. Its derivative is $\ln X/(X-1)$ and its limit as $X\to0^+$ is zero, which fixes the identity without numerical continuation.

